\documentclass[11pt]{article}
\usepackage[letterpaper,margin=0.82in]{geometry}
\usepackage[T1]{fontenc}
\usepackage{lmodern,amsmath,amssymb,amsthm,mathtools,booktabs,longtable,array}
\usepackage{microtype,xcolor,hyperref,xurl,fancyhdr,enumitem}
\hypersetup{colorlinks=true,linkcolor=blue!45!black,urlcolor=blue!50!black,citecolor=blue!45!black,pdftitle={Reusing witnesses in Diophantine consistency equations},pdfauthor={ChatGPT, prepared for Vladimir Reshetnikov}}
\setlength{\headheight}{14pt}
\pagestyle{fancy}\fancyhf{}\fancyhead[L]{\small Diophantine consistency: cofinal witnesses}\fancyhead[R]{\small 24 October 2026}\fancyfoot[C]{\thepage}
\setlength{\parskip}{4pt}\setlength{\parindent}{0pt}
\setlist{itemsep=3pt,topsep=4pt}
\newcommand{\Z}{\mathbb Z}\newcommand{\N}{\mathbb N}
\newcommand{\Con}{\operatorname{Con}}\newcommand{\ZFC}{\mathrm{ZFC}}
\newcommand{\Ccal}{C_{222}}\newcommand{\htop}{\operatorname{ht}}
\newcommand{\file}[1]{\nolinkurl{#1}}
\newtheorem{theorem}{Theorem}[section]
\newtheorem{lemma}[theorem]{Lemma}
\newtheorem{proposition}[theorem]{Proposition}
\newtheorem{corollary}[theorem]{Corollary}
\theoremstyle{definition}\newtheorem{definition}[theorem]{Definition}
\theoremstyle{remark}\newtheorem{remark}[theorem]{Remark}
\title{Reusing witnesses in Diophantine consistency equations\\[5pt]\large Deterministic CRT moduli, cofinal horizons, and shorter Pell transport}
\author{AI-assisted research prepared for Vladimir Reshetnikov\\\small ProveIt research program}
\date{24 October 2026}
\begin{document}
\maketitle
\begin{abstract}
This packet gives explicit ordinary integer polynomials representing finite
refutability of ZFC, with complete coefficient data and ordinary proofs of the
new transformations. The principal new idea is to reuse a power already
certified by a proof-certificate construction as the size bound for its
Chinese remainder encoding. A freely chosen encoding modulus can then be
replaced by a deterministic polynomial in existing coordinates. Independently,
the arbitrary padded proof horizon can be restricted to a cofinal family of
squares, and the sign of a proposed factorial output can be recovered from
existing power and ceiling equations. Choosing the modulus as a product of
consecutive integers makes its sign unconditional. The principal result is an explicit quartic
in 508 unrestricted integers with 6,545 monomials and coefficient height
996,521, improving the matched 517-variable parent in variables, residuals,
terms, and height. A separate positive-domain construction has ten unknowns,
degree 288, and an 89-operation evaluator. A simpler 512-variable checkpoint
and five degree/cost companions are retained. The direct construction
encodes the numerical proof calculus $\Ccal$; it does not simulate K215 or
instantiate Jones's universal polynomial. All claims of improvement are
relative to the frozen ProveIt sources identified below. No global minimum,
literature-wide priority, proof of consistency, or new end-to-end proof-assistant
verification is asserted.
\end{abstract}

\setcounter{tocdepth}{1}
\begingroup
\setlength{\parskip}{0pt}
\small
\tableofcontents
\endgroup
\newpage
\section{Results, conventions, and comparison boundary}
\label{sec:results}

We use $\N=\{0,1,2,\ldots\}$. Unless a result explicitly says
``strictly positive'', every unknown ranges independently over $\Z$.
A polynomial abbreviation is an instruction to substitute an integer
polynomial; it introduces no quantified unknown. Our equations contain no
variable exponent, factorial operation, division, square-root operation, or
unspecified coefficient in their final arithmetic representations.

The repository was frozen at commit
\begin{center}\small\file{45b399ec51dfbecb891dd51d8a9e4b6b7c361ac6}.\end{center}
Both the active candidates and the eight archives then present in
\file{docs/incoming} were examined. This matters because the overview of the
older manuscript collection is behind the most recent incoming constructions.
In particular, the matched direct parent is the 517-variable integer quartic
in the incoming Round 14 packet, rather than an older 569-, 604-, or
663-variable construction \cite{round14,repo}.

\begin{center}
\begin{tabular}{lrrrr}
\toprule
Integer quartic & Variables & Rows & Terms & Height\\
\midrule
Matched incoming parent & 517 & 310 & 6,774 & 3,904,267\\
Cofinal/modulus checkpoint & 512 & 310 & 6,742 & 3,904,268\\
\textbf{Signed-factorial principal} & \textbf{508} & \textbf{308}
 & \textbf{6,545} & \textbf{996,521}\\
\addlinespace
Incoming lower-height alternative & 519 & 321 & 6,253 & 411,060\\
\bottomrule
\end{tabular}
\end{center}
All three have exact degree four. The 512-variable checkpoint slightly
increases height; the principal 508-variable equation improves every
displayed quantity relative to the matched parent. A separate family has
ten strictly positive unknowns, with operation/degree pairs
$89/288$, $91/283$, $92/278$, $94/262$, $96/258$, and $97/254$.
Section~\ref{sec:positive} specifies those polynomials and their domains.
The separate 519-variable incoming equation retains fewer terms and smaller
coefficients than the 508-variable result; it remains a useful tradeoff.


Here ``rows'' means integral quadratic polynomials $R_j$ in the representation
$P=\sum_j R_j^2$, and coefficient height means the largest absolute coefficient
after the full sum has been expanded and like monomials collected. Arithmetic
operation counts for the compact equations count each live binary addition,
subtraction, and multiplication, including squares and multiplication by a
fixed constant. They do not measure bit complexity: those equations have very
large fixed native coefficients.

There is no single established meaning of ``simplest'' that orders all these
equations. A few unknowns, low degree, small coefficients, few terms, short
evaluation circuits, and short correctness proofs are different objectives.
The direct quartics improve a modest-coefficient, independently encoded branch
of the repository. They do not supersede the repository's separate
37-integer quartic with large native coefficients, its ten-integer
higher-degree constructions, or all of its positive-domain forms.

\subsection{What is proved in ZFC}
Fix the finite syntactic ZFC-refutation relation of the retained source
calibration and write
\[
 \Con(\ZFC)\;:\Longleftrightarrow\;
 \neg\exists p\in\N\;\operatorname{Prf}_{\ZFC}(p,\bot).
\]
For each asserted equation the target is the implication in both directions,
\[
 \ZFC\vdash
 \left(\exists\mathbf z\in D^n\;P(\mathbf z)=0\right)
 \longleftrightarrow
 \left(\exists p\in\N\;\operatorname{Prf}_{\ZFC}(p,\bot)\right),
 \tag{1}\label{eq:target}
\]
where $D$ is the stated domain. Negation gives the requested consistency
equivalence. The proof is an ordinary mathematical proof formalizable in ZFC,
composed with precisely identified inherited calculus and arithmetic theorems.
No truth or soundness assumption about ZFC is used to infer a refutation from a
polynomial zero. Conversely, a finite syntactic refutation is used as finite
data for constructing integer witnesses.

The artifact checks establish literal polynomial identities, coordinate maps,
coefficient expansions, and numerical metrics. They do not establish
unsolvability, decide consistency, or replace the unbounded arithmetic proofs.
The supplied Lean interface files document the inherited calibration; their
entire transitive import closure has not been rebuilt for this packet.

\section{The direct proof-certificate interface}
\label{sec:interface}

The direct construction begins with the repository's numerical calculus
$\Ccal$. Its implication, membership, universal-quantifier, and equality
constructors are fixed integer polynomials in the schema parameters; it has
finitely many axiom and inference branches, a truth code $0$, and a falsity
code $1$. The complete constructor formulas and the guarded axiom condition
are retained with the source documents. In the square-shell variant used
here, the corresponding syntax is related to the original coding by the
retained finite syntax-transfer argument; it is not an arbitrary relabeling.

The inherited logical calibration is
\[
 \Ccal\text{ derives }1
 \quad\Longleftrightarrow\quad
 \mathrm{ZF}\text{ has a finite refutation}
 \quad\Longleftrightarrow\quad
 \ZFC\text{ has a finite refutation}.
 \tag{2}\label{eq:calibration}
\]
The named interfaces are
\file{KEndpoint.c222_iff_zf_refutation} and
\file{ConstructibleChoice.zf_refutation_iff_zfc_refutation}.
The second uses the constructible-universe relative consistency argument.
The filename of the first interface does not make the present arithmetic
construction a simulation of a Turing machine.

\subsection{The parent theorem used by the new reduction}
The retained direct theorem has the following concrete properties.

\begin{enumerate}
\item A finite refutation can be expanded to a finite proof tree. After
choosing any sufficiently large integer horizon $t\ge16$, truth steps pad
the backward proof expansion to exactly $t$ rows. The lower bound on $t$
depends only on finitely many proof-tree parameters and constructor labels,
including the fixed truth-padding tuple.
\item The semantic factorial is $F=t!$, and the label of code $j$ is assigned
the positive weight $W_j=1+jF$ for $0\le j\le t$.
\item A backward expansion uses positive accumulators $S_i$ and $S_{i+1}$,
which are products of the weights of the currently pending formulas. A row
with conclusion $v$ and premises $p,q$ satisfies
\[
 S_{i+1}W_v=S_iW_pW_q.
 \tag{3}\label{eq:flow}
\]
Unused premises have code $0$, whose weight is $1$. The endpoints are
$S_0=F+1$ and $S_t=1$.
\item The first CRT encoding uses moduli $1+(i+1)y$, $0\le i\le t$,
where $y$ is any sufficiently large positive multiple of $F$. In particular,
$S_i\le y$ and $y\ge F+1$ suffice. The raw endpoint slack is
$r_0=y-F-1\ge0$.
\item The factorial module already certifies the power
\[
 \rho=Ft^2,\qquad V=(\rho+1)^t.
 \tag{4}\label{eq:oldpower}
\]
Its noncircular signed-power and digit argument recovers $F=t!$ and the
intended value of $V$ from every parent zero.
\item A second, prime-power CRT encoding and shared Pell construction
compile the bounded local conditions into a finite integral quadratic
system. The integer version imposes a finite selected list of nonnegativity
conditions with private three-variable representations. Exact monic
elimination and integer translation produce the delivered 517-variable
quartic.
\end{enumerate}

These are the actual completeness and soundness properties proved in
\file{FACTORIAL_CARRIER_PROOF.md},
\file{DERIVED_LABEL_RANGE_PROOF.md},
\file{HYBRID_PROOF.md}, and the retained signed-source proof. They are
included in the packet, together with the literal parent equation and its
source ledger. The new argument does not replace any of these statements by
an asymptotic assertion about an unspecified universal polynomial.

\subsection{Why the multiplicative flow is sound}
The positive weights $W_1,\ldots,W_t$ are pairwise coprime. Indeed, a prime
dividing both $W_j$ and $W_k$ would divide $(j-k)F$ while not dividing $F$.
It would therefore be at most $t$, hence would divide $F=t!$, a contradiction.
Multiplying \eqref{eq:flow}, cancelling the overlapping positive accumulators,
and taking a prime valuation private to each $W_j$ proves
\[
 \#\{\text{conclusions with label }j\}
 -\#\{\text{premises with label }j\}
 =\begin{cases}1&j=1,\\0&j>1.\end{cases}
 \tag{5}\label{eq:balance}
\]
Cycles in the rule multiset cause no gap. Let $D$ be the closure of truth
under its finitely many valid rule instances. If falsity were outside $D$,
each rule concluding outside $D$ would have a premise outside $D$.
The number of outside conclusions would be at most the number of outside
premises, contradicting the sum of \eqref{eq:balance} over labels outside $D$.
Thus falsity has a finite derivation. This is the logical reason the compact
unordered certificate still represents refutability.

\section{Cofinal witness restriction}
\label{sec:cofinal}

The new restriction is applied to a freely chosen part of a completeness
construction. It is not a surjectivity claim for squares.

\begin{lemma}[Cofinal selection]
Suppose that every finite certificate $c$ has realizations with every integer
parameter $t\ge L(c)$, with all later witnesses chosen after $t$. Let
$\phi:\Z\to\N$ have unbounded image and let the extra equation
$t=\phi(u)$ be expressible by the permitted polynomial system. If the
original system is sound for every such $t$, the restricted system has a
zero if and only if a finite certificate exists.
\end{lemma}
\begin{proof}
A restricted zero is an original zero. Conversely choose $u$ with
$\phi(u)\ge L(c)$ and then carry out the original completeness construction
at that parameter. The order of choices is essential: witnesses depending on
$t$ are regenerated after $u$ has been selected.
\end{proof}

For the present certificate, take
\[
 t=16+u^2. \tag{6}\label{eq:squarehorizon}
\]
There is no congruence condition on the padded proof horizon. Given the finite
proof tree, choose $|u|$ large enough to dominate its step count, every local
schema parameter, and every complete constructor-label bound. Padding by
truth remains available. The factorial, weights, accumulators, both CRT
encodings, and Pell witnesses are then constructed afresh.

The original unrestricted-integer implementation used
\[
 t-16=a^2+b^2+c(c+1).
\]
Replacing it by \eqref{eq:squarehorizon} removes two private coordinates and
keeps the row quadratic. The classical equivalence
\[
 n\ge0\quad\Longleftrightarrow\quad
 \exists a,b,c\in\Z\;n=a^2+b^2+c(c+1)
 \tag{7}\label{eq:three}
\]
follows from the three-square theorem applied to $4n+1$; see also the explicit
formulation in \cite[Lemma 6.1]{bayer}. The reverse implication follows
directly from $c(c+1)\ge0$ for integral $c$.

\begin{remark}[The restriction is not valid for an arbitrary coordinate]
An equation forcing a natural witness to equal $2$ has a solution, whereas
forcing that witness to be a square destroys all its solutions. Similarly,
an unbounded admissible set such as $\{4k+2:k\in\N\}$ contains no square.
Unboundedness alone is insufficient. We use the parent's stronger fact that
every sufficiently large horizon is admissible, with the other witnesses
chosen afterwards.
\end{remark}

\section{A deterministic CRT modulus from an existing power}
\label{sec:modulus}

\begin{lemma}[Pending-product bound]
Let a proof tree be expanded backwards for at most $t$ steps, with each
step replacing one pending conclusion by at most two premises. Suppose all
labels lie in $[0,t]$, $t\ge1$, and $F\ge1$. Put
$A=1+Ft^2$ and $V=A^t$. Every accumulator in the expansion, including
truth padding, is at most $V^2$.
\end{lemma}
\begin{proof}
The initial pending multiset has one member. After $s$ steps its cardinality
is at most $s+1$, hence at most $t+1$. Each factor is
$W_j=1+jF\le1+tF\le A$. Consequently
\[
 S_i\le A^{t+1}\le A^{2t}=V^2.
\]
Truth has weight $1$ and leaves this bound unchanged.
\end{proof}

The expression $V^2$ already belongs to the same arithmetic certificate:
the parent does not need a new exponentiation request to obtain it. Moreover,
\[
 V=(1+Ft^2)^t\equiv1\pmod F.
 \tag{8}\label{eq:Vcong}
\]
This congruence permits a deterministic modulus with both a size property
and the necessary divisibility property.

\begin{proposition}[A family of polynomial moduli]
Let integers $\kappa\ge2$ and $d\ge1$ satisfy $(\kappa-1)d\ge2$.
Replace the freely chosen modulus and its slack by
\[
 F\ge d,\qquad y=V^2+\kappa F-1.
 \tag{9}\label{eq:modfamily}
\]
For arbitrary signed $V$, these conditions imply
$r_0=y-F-1=V^2+(\kappa-1)F-2\ge0$.
For canonical $F=t!$ and $V$ from \eqref{eq:oldpower}, they also imply
$F\mid y$ and $y>V^2$. Thus they supply every first-CRT condition required
by the parent completeness construction.
\end{proposition}
\begin{proof}
The slack inequality is immediate from the displayed formula and the
condition on $d$. Congruence \eqref{eq:Vcong} gives
$V^2+\kappa F-1\equiv0\pmod F$. Finally
$\kappa F-1\ge1$, so the pending-product lemma implies every required
accumulator is strictly smaller than $y$.
\end{proof}

The delivered 512-variable equation uses $\kappa=3,d=1$:
\[
 \boxed{\quad y=V^2+3F-1,\qquad F\ge1.\quad}
 \tag{10}\label{eq:chosenmod}
\]
Strengthening the existing $F\ge0$ row to $F\ge1$ changes a constant
but introduces no coordinate. The former three private witnesses for
$r_0\ge0$ are deleted. The new modulus equation is a single quadratic
row. There is no requirement to store $V^2$ as another unknown.

\subsection{The CRT construction remains valid}
For completeness, $F=t!\mid y$. Suppose a prime $p$ divides both
$m_i=1+(i+1)y$ and $m_j=1+(j+1)y$, where $0\le i<j\le t$.
It does not divide $y$, so $p\mid j-i$ and $p\le t$.
But then $p\mid F\mid y$, a contradiction. Therefore these moduli are
pairwise coprime. The Chinese remainder theorem supplies a nonnegative $x$
with $x\equiv S_i\pmod {m_i}$ for all $i$. As
$0<S_i\le y<m_i$, the corresponding Euclidean quotients are nonnegative.
The initial and final residues are exactly $F+1$ and $1$.

The order of the construction is
\[
 \text{finite proof}\ \longrightarrow\ t\ \longrightarrow\ F,V,y
 \ \longrightarrow\ x\ \longrightarrow\text{second CRT and Pell witnesses}.
\]
In particular, the deterministic choice of $y$ has no dependence on the
later code $x$. The second CRT construction still chooses its representatives
at the original strengthened prime-power precision. No claim that one fixed
small representative meets all amplified congruences is introduced.

\section{The complete 512-variable integer quartic}
\label{sec:512}

Let $R_0,\ldots,R_{309}$ be the literal rows in
\file{direct/equation512.json}, with variables
$x_0,\ldots,x_{511}$ in the recorded order. Define
\[
 \boxed{P_{512}(x_0,\ldots,x_{511})
       =\sum_{j=0}^{309}R_j(x_0,\ldots,x_{511})^2.}
 \tag{11}\label{eq:P512}
\]
This defines one fixed integer polynomial, not a family indexed by a program
or a theorem predicate. The complete collected coefficient list is
\file{direct/polynomial512.json}.

The parent is already in translated coordinates. Its semantic horizon is
$t=15+x_0$, not $16+x_0$; the old nonnegative horizon slack is $x_0-1$.
The new horizon row is therefore $x_0-1-u^2=0$, which is precisely
\eqref{eq:squarehorizon}. The new rows before removing unused coordinates
are recorded verbatim in \file{direct/specialization_certificate.json}.
The exact old coordinate indices removed are
\[
 329,\ 330,\ 343,\ 344,\ 345.
\]
Every other row is unchanged apart from the final order-preserving renumbering.
No coordinate whose value is used elsewhere is deleted.

\begin{theorem}
The polynomial $P_{512}$ has exact degree four, 6,742 nonzero collected
monomials, coefficient height 3,904,268, and exactly 512 occurring variables.
ZFC proves
\[
 \bigl[\forall\mathbf x\in\Z^{512}\;P_{512}(\mathbf x)\ne0\bigr]
 \longleftrightarrow\Con(\ZFC).
\]
\end{theorem}
\begin{proof}
If $P_{512}=0$, each integer square in \eqref{eq:P512} is zero. The
new horizon row implies the old horizon inequality, the strengthened
factorial row implies $F\ge1$, and \eqref{eq:chosenmod} gives
$r_0=V^2+2F-2\ge0$ without knowing the meaning of $V$. By
\eqref{eq:three}, the deleted private nonnegativity witnesses can be
restored. The unchanged rows and the restored witnesses give a zero of the
parent 517-variable system. The retained parent soundness theorem and
\eqref{eq:calibration} yield a finite ZFC refutation.

Conversely, translate a finite ZFC refutation into a finite $\Ccal$
derivation using the inherited calibration. Expand it to a finite proof tree.
Choose the cofinal horizon \eqref{eq:squarehorizon} beyond all finite bounds,
then set $F=t!$ and $V=(Ft^2+1)^t$. The preceding section constructs the
deterministic admissible modulus $y$ and the first CRT code. The parent's
remaining constructive completeness proof supplies the second CRT data and
all shared Pell witnesses. Since $F-1\ge0$, the strengthened factorial
guard has three-square witnesses. These values satisfy every new row. Exact
parent projection and translation then give a zero of $P_{512}$.

The coefficient and row data establish the stated finite metrics by exact
expansion. Degree four also follows intrinsically: a nonzero homogeneous
quadratic part occurs among the residuals, and the sum of their squares
cannot vanish identically over $\mathbb R$. The quantified transformations
use finite syntax, induction, integer arithmetic, CRT, the three-square
theorem, and the retained Pell/calibration results, all formalizable in ZFC.
Negating the zero-existence equivalence gives the displayed statement.
\end{proof}

This is an equivalence of existence. An arbitrary old zero with a nonsquare
horizon slack need not be a new zero. The converse rebuilds a certificate
from the represented finite refutation. Confusing these two kinds of
reduction would invalidate the completeness argument.

\section{A signed factorial output: the 508-variable refinement}
\label{sec:508}

The preceding 512-variable equation is a useful control because it changes
only three rows of the final parent. A further change at the arithmetic
source recovers the factorial output's sign from the equations. Its three
private sign witnesses can therefore be removed. Sharing the existing square
also permits one additional exact projection in the rebuilt system.

\subsection{The changed source}
Keep $F$ as an unrestricted integer, retain the cofinal horizon
$t=16+u^2$, and change the arithmetic definitions to
\[
 \rho=F^2t^2,\qquad
 R=2+t+x+y+F^2,\qquad K=2R,\qquad
 \boxed{y=V(V-1).}
 \tag{12}\label{eq:signedchanges}
\]
The power request remains $V=(\rho+1)^t$, but its meaning is established
only after the initial sign argument. The existing source already has the
wire for $F^2$ because its reciprocal product uses
$B=\binom{F^2+t}{t}$. Thus \eqref{eq:signedchanges} adds no fresh square
coordinate. The formula for $y$ is quadratic in the existing $V,y$.

Deleting the factorial guard is not justified by silently assuming $F=t!$.
Indeed, before the power requests are interpreted, $F$ and $V$ may have
either sign. The crucial simplification is that every integer $v$ satisfies
$v(v-1)\ge0$: either $v\ge1$ or $v\le0$. Consequently the literal new
modulus row gives
\[
 y=V(V-1)\ge0,\qquad y+1\ge1
 \tag{13}\label{eq:consecutivepositive}
\]
before any indexed-Pell or power interpretation. This use of discreteness
does not hold for arbitrary real $V$; the domain here is $\Z$.
All independent seed indices are therefore positive, in particular
\[
 t,\ K,\ F^2t^2+t,\ F^2+t,\ y+1,\ x+1\;>0.
 \tag{14}\label{eq:earlyy}
\]
The remaining seeds are the retained positive encoded representatives.
The local constructor bound is also positive on those
representatives. Here $x\ge0$ and every unchanged seed guard are retained.
In particular $R\ge2+t>0$ after \eqref{eq:earlyy}.

The unchanged common budget and guarded frame norm give the same parameter
bounds as in the parent. Every seed inequality implies a positive Pell ordinate. Their shared
auxiliary product is nonzero, so its norm establishes the previously
required auxiliary-coordinate bound. The signed indexed-Pell theorem is
therefore licensed in the same manner as in the parent. The old positivity
certificate is replaced by these actual seed formulas; it is not cited
unchanged after the deletion of its factorial premise. No exceptional
zero-output case or parity restriction on the horizon is needed.

\subsection{Recovering powers and the factorial in the correct order}
Let $D$ denote the proposed factorial binomial output. The retained natural
remainders give $D\ge1$ and a common radix $q\ge2$. The request $q^t$
orients the shared $t$-indexed Pell pair. Thus
\[
 V=(F^2t^2+1)^t\ge1.
\]
The raw ceiling rows give $DF=V+s$, where $s\ge0$. Since $D\ge1$,
this forces $F\ge1$ before any factorial value is assumed. Since $t\ge16$,
$V\ge F^2t^2+1\ge F+2$, and therefore
\[
 y-F-1=V(V-1)-F-1\ge0.
\]
The old endpoint slack can now be recovered.

The reciprocal raw index row is $C+F^2+1=B i_q$, with $C\ge1$ and
$B\ge0$ still guarded. It gives $B>0$. In the signed power request
$N=B^K$, the only possible inverse case is $B=1,N=-1$. But then the
other base at exponent $K$ is $N+V+1\ge2$. That request orients the
$K$-indexed pair and excludes the inverse case. The common radix has its
intended value
\[
 q=N(N+V+1)^K>V^2.
\]
Every remaining seed is similarly oriented by a request at base $q$ or
$q+1$. This establishes all power meanings without first interpreting the
factorial or the binomials whose digits are under validation.

Now $\rho=F^2t^2\ge t$, and for $0\le j\le t$,
\[
 \binom{\rho+t}{j}\le\binom{\rho+t}{t}
 \le(\rho+1)^t=V<q.
\]
The first digit graph therefore really gives $D=\binom{\rho+t}{t}$.
The exact factorial argument is recorded next, rather than inferred from
finite sample calculations.

\begin{lemma}[Self-bounding factorial at the squared radius]
Let $F,t$ be positive integers, $\rho=F^2t^2$,
$V=(\rho+1)^t$, and $D=\binom{\rho+t}{t}$. Then
\[
 F-1<V/D\le F\quad\Longleftrightarrow\quad F=t!.
\]
\end{lemma}
\begin{proof}
Write $G=t!$. The exact identity is
\[
 \frac VD=G\prod_{i=1}^t
   \left(1-\frac{i-1}{\rho+i}\right).
\]
All factors lie in $(0,1]$, and
\[
 \sum_{i=1}^t\frac{i-1}{\rho+i}
 <\frac{1}{2F^2}.
\]
The elementary product bound gives
$G(1-1/(2F^2))<V/D\le G$.
If $F-1<V/D\le F$, then $G>F-1$, so $G\ge F$.
On the other hand,
\[
 G<\frac{F}{1-1/(2F^2)}
   =F+\frac{F}{2F^2-1}\le F+1,
\]
and integrality gives $G\le F$. Thus $F=G$.
Conversely, if $F=G$, then
$F-1<F-1/(2F)<V/D\le F$; when $F=1$ the same strict inequality
also holds. The ceiling remainder $DF-V$ lies in $[0,D-1]$, giving
the required natural ceiling witnesses.
\end{proof}

The source ceiling rows are precisely the inequality in this lemma, so
$F=t!$. The second digit is also safe: for $0\le j\le t$,
\[
 \binom{F^2+t}{j}\le(F^2+t)^t\le(\rho+1)^{2t}=V^2<q.
\]
It recovers the reciprocal product $B$ and $B\ge2$. The later amplified
CRT and derived-label bounds are preserved: after $F\ge1$ the new
$R=2+t+x+y+F^2$ is at least the old $2+t+x+y+F$. The same local
majorant exponent, factorial growth floor, representative bounds, and
strengthened CRT precision therefore apply. No preceding bootstrap step
used that later factorial growth estimate.

\subsection{Completeness and the literal polynomial}
Choose the cofinal padded horizon as before and set $F=t!$,
$\rho=F^2t^2$, and $V=(\rho+1)^t$. A slightly stronger pending-product
bound is useful. Put $a=Ft\ge16$, $L=a+1$, and $B_0=a^2+1$.
There are at most $t+1$ pending weights, each at most $L$, so
$S_i\le L^{t+1}$. Since $a\ge3$ and $t\ge2$,
\[
 B_0>(a-1)L,\qquad (a-1)^t\ge(a-1)^2\ge L.
\]
It follows, using integer powers only, that
\[
 V=B_0^t>L^t(a-1)^t\ge L^{t+1}\ge S_i.
\]
Now $V\ge3$, so the prescribed modulus $y=V(V-1)$ is at least $V$
and exceeds every accumulator. Also $F\mid(V-1)$ because
$V\equiv1\pmod F$, and hence $F\mid y$. We have already established
$y\ge F+1$. These are exactly the inherited CRT requirements; the earlier
sufficient convenience $y>V^2$ is unnecessary. The first CRT is valid. The larger new
$R$ satisfies every subsequent lower bound. The factorial lemma supplies
the ceiling witnesses, the unchanged digit arguments supply their quotients,
and simultaneous Pell completeness supplies the common frame. Thus every
finite refutation gives a zero of the new arithmetic system.

The source is then factored, projected through monic integral definitions,
and translated by the same recorded pipeline as its matched parent. The
principal result is
\[
 \boxed{P_{508}(x_0,\ldots,x_{507})
       =\sum_{j=0}^{307}Q_j(x_0,\ldots,x_{507})^2=0.}
 \tag{15}\label{eq:P508}
\]
The full $Q_j$ are printed in Appendix~\ref{app:principal-residuals} and supplied in
\file{direct/principal/equation.json}. Their full collected expansion is
\file{direct/principal/expanded_polynomial.json}.

\begin{theorem}[Principal unrestricted-integer equation]
The polynomial in \eqref{eq:P508} has 508 occurring integer variables,
308 quadratic residuals, exact degree four, 6,545 collected monomials, and
coefficient height 996,521. ZFC proves
\[
 \bigl[\forall\mathbf x\in\Z^{508}\;P_{508}(\mathbf x)\ne0\bigr]
 \longleftrightarrow\Con(\ZFC).
\]
\end{theorem}
\begin{proof}
The product of consecutive integers establishes the modulus sign, and hence
all seed signs, without assuming a factorial or power meaning. The ordered recovery gives the
same bounded local proof-certificate semantics, whose soundness gives
\eqref{eq:calibration}. The constructive paragraph above proves the
converse from any finite refutation. Every source modification is
identified by an exact quadratic-row certificate. The subsequent monic
eliminations and integer translations have recorded integral inverse maps,
so they preserve integer-zero existence. Finally a sum of integral squares
vanishes exactly when all its residuals vanish. Exact coefficient expansion
establishes the finite metrics, including a nonzero degree-four part.
All quantified mathematical steps are formalizable in ZFC, yielding
\eqref{eq:target} for $D=\Z$ and $n=508$.
\end{proof}

The matched parent has 517 variables, 310 rows, 6,774 terms, and coefficient
height 3,904,267. This refinement improves all four of those quantities
while retaining degree four and unrestricted integer unknowns. It removes
about $1.74\%$ of the variables and $3.38\%$ of the monomials; coefficient
height falls by about $74.48\%$. Those are measured improvements in this
particular direct construction, not a claim that 508 is the smallest
possible integer-variable count.

\section{Ten positive unknowns with an 89-operation evaluator}
\label{sec:positive}

The direct result above keeps modest coefficients and unrestricted integer
unknowns. A complementary calculation shortens the existing compact native
route. This section's domain is explicitly $\Z_{>0}^{10}$, and its large
fixed coefficients are inherited from the coherent compact44 BFS tuple.
It is not a 10-variable unrestricted-integer result.

\begin{center}
\begin{tabular}{lrrr}
\toprule
Polynomial & Variables & Operations & Exact degree\\
\midrule
Uncentered & 10 & 89 & 288\\
Upper center & 10 & 91 & 283\\
Sharp center & 10 & 92 & 278\\
Direct, uncentered & 10 & 94 & 262\\
Direct, upper center & 10 & 96 & 258\\
Direct, sharp center & 10 & 97 & 254\\
\bottomrule
\end{tabular}
\end{center}

The centered parents in the incoming saturated-Pell packet cost, in the
same order, $92/283$, $93/278$, $98/258$, and $99/254$, where each pair is
operations/degree \cite{saturated}. Thus all four centered companions
improve cost at unchanged degree. The uncentered 89-operation equation is
a new cost/degree tradeoff. The earlier square-gap and auxiliary-index
ideas in the incoming structural-compression packet are explicitly credited
\cite{structural}; the new step is their combination with coprime Pell
fusion and the resulting replacement of the transport majorant.

\subsection{Complete polynomial definition}
Use the positive unknowns
\[
 J,F,\alpha,i,T,w,Y,n,\eta,\zeta.
\]
On the direct chart replace $w$ by a positive unknown $x$. The four
constants $\beta,K_0,M_C,M_F$ are the entire compact44 BFS tuple supplied
in \file{positive/coefficients/}. In particular $\beta=2^{590069}-1$.
The full fixed finite recipe, including its 436-character seed and all
finite index bounds, is \file{positive/coefficients/recipe_compact44.py}.
It reconstructs every coefficient and checks its complete raw-byte hash.
The high mask $M_F$ is already shifted; no second shift or mixture with
the differently ordered original44 tuple is permitted.

Define only polynomial abbreviations:
\begin{align*}
 q&=\beta J+1,& \mu&=q^2+1,\\
 k&=\eta+\zeta,&c&=kY+\eta,\\
 X&=wq\quad\text{(ordinary)},&&
 X=q+x\quad\text{(direct)},\\
 A&=\mu Y(X+1)+1,&H&=4A-5,&\Delta&=A^2-1,\\
 V_0&=X(\mu Y)^2,&E&=V_0(V_0+1)k^2+1,\\
 R&=\bigl[(q-1)(q-F)+\alpha\bigr](q^2-1)\\
  &\qquad +(M_C+qM_F)J,\\
 e&=k-R-1,&C&=q-F-\alpha,&U&=ic^2,\\
 D_f&=1+\Delta U^2E,&Q&=\Delta^2U^2E,&\delta&=Q-1,\\
 b&=c(T+\alpha)-R,&V&=D_fb-c,&D_y&=QV^2-1,\\
 K_+&=X(3c-2X)+2.
\end{align*}
The ordinary chart uses
\[
 B=(K_0+w)C-F,\quad h=q-1,\quad P=q^3(q-1),\quad D_{\rm low}=0,
\]
and the direct chart uses
\[
 B=(qK_0+X)C-qF,\quad h=q(q-1),\quad
 P=q^4(q-1),\quad D_{\rm low}=(q-1)X.
\]
Select a center
\[
 C_g=0,\qquad C_g=(A+1)UE,\qquad\text{or}\qquad C_g=AUE,
\]
respectively. Finally set
\begin{align*}
 M&=P\mu^2H,\\
 Z&=P\mu^2\{H(n+V+C_g-bB)+K_+\}\\
  &\qquad +H\{\mu^2(hY-q^3B-D_{\rm low})-Pe\},\\
 s&=\delta Z^2,\qquad u=\delta M^2ED_f,\qquad v=M^2D_y.
\end{align*}
The single equation, for any of the six choices, is
\[
 \boxed{N=(s-u-v)^2-4uv=0.}
 \tag{16}\label{eq:positiveN}
\]
All abbreviations use only addition, subtraction, multiplication, and fixed
nonnegative integer exponents. The literal JSON circuits expand every one
of them; no square root or auxiliary predicate occurs in $N$.

\subsection{The auxiliary coordinate is a sufficient transport bound}
The inherited coprime recovery, before any transport sign conclusion, gives
positive integers $\tau,f,y$ with
\[
 E=\tau^2,\quad D_f=f^2,\quad I=U\tau,\quad
 Q=(\Delta I)^2,\quad QV^2-(Q-1)y^2=1.
 \tag{17}\label{eq:recovered}
\]
For clarity, here $I$ is a proof-only recovered integer, not an additional
unknown. The coprime step is valid because $D_f\equiv1\pmod E$.
The two rational radicands $ED_f$ and $D_y/\delta$ cannot coincide:
after multiplying by $\delta$ and reducing modulo $E$, equality would give
$0\equiv-1$. Thus the rational signed-sum argument makes both roots
rational; square-class coprimality separates $E$ and $D_f$, and the
retained adjacent-coefficient descent makes $y$ integral. This order does
not presuppose the later projection or the sign of $B$.

Write $\psi_A(m)$ for the Pell ordinate at base $A$.
The retained uniform folding and first-index arguments give
$c=\psi_A(p)$ for an odd positive $p$, together with
$c>8q^5$, $c>2X$, $H>3X$, and $K_+>0$.
The native mask bound
$R<q^4+(\alpha+2)q^2<c\alpha/4$ then gives
\[
 b=c(T+\alpha)-R>c.
\]
Since $c\ge3$, the odd index $p$ is at least three, so
$c\ge4A^2-1=4\Delta+3$.

\begin{lemma}[Fifth-index gap after fusion]
Under the recovered conditions above,
\[
 b>8Q,\qquad g:=\tau f<Q,\qquad 0<y-V<b/4.
\]
\end{lemma}
\begin{proof}
From $b>c$,
\[
 V=D_fb-c>(D_f-1)c=\Delta I^2c>4Q.
\]
Put $S=\Delta I$, so $Q=S^2$. Pell classification applied to
$(SV)^2-(S^2-1)y^2=1$ gives $SV=\chi_S(\nu)$ with positive odd
$\nu$. The values at $\nu=1,3$ give $V=1,4Q-3$, both excluded.
Thus $\nu\ge5$ and
\[
 V\ge16Q^2-20Q+5>8Q^2.
\]
Since $D_f<Q$, this yields $b=(V+c)/D_f>8Q$.
Also $f<AI$, so $g<AUE<Q$.
Finally $Q-1>2D_f$ and
\[
 0<y-V=\frac{V^2-1}{(Q-1)(y+V)}
       <\frac{V}{2(Q-1)}<\frac b4.
\]
Each inequality precedes the transport sign recovery.
\end{proof}

For explanation only, divide the carrier by $M$ and put $W=1+hb$.
The exact identity is
\[
 \frac ZM=n+V+C_g+\frac{K_+}{H}+\frac{Y}{q^3}
          -\frac{WB}{h}-\frac e{\mu^2}-d,
 \tag{18}\label{eq:carrierpositive}
\]
where $d=0$ ordinarily and $d=X/q^4$ directly. At a norm zero,
$Z/M=\varepsilon_g g+\varepsilon_y y$.
Since $Q>X$, we have $g+d<2Q<b/4$. Also
$b-e=c(T+\alpha)-k+1>0$, so $e/\mu^2<b/4$ even when $e$ is negative.
If $B\le-1$, solving \eqref{eq:carrierpositive} for $n$ gives
\[
 n<\frac{3b}{4}+\frac{WB}{h}
   \le\frac{3b}{4}-b-\frac1h<0.
\]
This contradicts the domain of $n$. A positive center $C_g$ only decreases
the right side. Thus all six equations force $B\ge0$ using the existing
coordinate $b$ itself; the old separately computed transport majorant is
unnecessary. The cheapest pair can also discard the center.

\subsection{Completeness and the consistency theorem}
After $B\ge0$ is established, the unchanged source arguments recover
$C>0$, $3q<R<q^4-2$, the exact index $p=R$, denominator divisibilities,
and $X=2^R$. They reconstruct the same unscaled native source witness;
the scaled strong-Pell witness is not divided and reused without proof.
The retained native/controller/refutation calibration therefore applies.

Conversely, start with the parent's canonical data supplied by a finite
refutation. Choose a fresh strong Pell period with
$c^2\tau\mid I$, which is possible by the finite order of the
determinant-one Pell matrix modulo $c^2\tau$. Put
$i=I/(c^2\tau)>0$. The retained auxiliary-completeness construction gives
positive $T$, the auxiliary value $V$, and a gap $u_*=y-V>4I^2$.
It also gives
\[
 \kappa=K_+/H<c<I,\quad s_Y=Y/q^3<Y<c<I,\quad
 UE=I\tau<I^2.
\]
First work on the ordinary chart, with
$B_t=(K_0+w)C-F$ and $W=1+(q-1)b$.
The inherited canonical divisibilities make $\kappa$, $s_Y$,
$t_*=B_t/(q-1)\ge0$, and $r_*=e/\mu^2\ge0$ integers. Choose
\[
 n=g-C_g+u_*-\kappa-s_Y+Wt_*+r_*.
\]
For the uncentered choice $g>2c>\kappa+s_Y$ proves $n>0$.
For the centered choices, $0<C_g-g<2UE$ and the total negative loss
is less than $2I^2+2I<4I^2<u_*$. The norm then vanishes.
The direct chart takes $x=X-q>0$ and adds the positive integral correction
$(W_{\rm direct}-W_{\rm ordinary})t_*+X/q^4$ to $n$.
Here the inherited canonical power divisibility gives $q^4\mid X$.
Thus every finite refutation supplies positive zeros in both charts.

It follows that, for every polynomial $N$ in the table,
\[
 \ZFC\vdash
 [\forall\mathbf z\in\Z_{>0}^{10}\;N(\mathbf z)\ne0]
 \longleftrightarrow\Con(\ZFC).
\]
The complete new proof, with the precise inherited theorem interfaces,
is \file{positive/PROOF.md}.

\subsection{Exact degree, rather than a syntactic upper estimate}
Let $q_0=\beta J$, $k_0=\eta+\zeta$, and $x_0=q_0+x$.
Uniform formal scaling of all ten variables gives the following nonzero
highest homogeneous parts for the uncentered pair:
\begin{align*}
 N_{288,\mathrm{top}}
 &=-2^{10}q_0^{152}w^{30}Y^{62}i^{10}k_0^{32}(q_0-F)^2,\\
 N_{262,\mathrm{top}}
 &=-2^{10}q_0^{126}x_0^{30}Y^{62}i^{10}k_0^{32}(q_0-F)^2.
\end{align*}
They are nonzero for the actual positive fixed $\beta$.
In the formal Laurent-series field, the rational remainder in
$Z/M=V+C_g+\mathcal E$ has degree at most six. The centered radical
defects have degrees at least seventeen, so the changed transport cannot
alter their leading terms. Pairing the two signs of each radical yields
the exact centered degrees in the table. The proof provides the complete
leading forms and global degree bound; exact one-variable specialization
checks are supplementary transcription controls, not the sole degree proof.


\section{Verification, reproducibility, and evidence boundaries}
\label{sec:verify}
The release contains a portable \file{verify_all.py}. It checks the exact
finite artifacts without querying a service or fetching a repository. Python
3.10 or later and the standard library suffice for the supported checkers.
The complete equation specifications are the literal JSON files, not the
human-readable metric summaries.

The principal direct checker independently compares the modified source
against its frozen parent, checks the complete square budget and the new
seed/sign prerequisite rows, and replays the factorization, monic projection,
and translation maps. It expands every quadratic residual square and compares
every collected integer coefficient with the supplied expansion. It does
not import the equation builder to decide whether the final coefficients are
correct. The 512-variable checker separately checks all three changed rows
and the exact order-preserving coordinate deletion.

For the positive companions, the independent verifier compares all six
complete parent-to-child graphs, reconstructs all four actual fixed
coefficients, checks every live operation, expands the universal two-root
norm identity, and checks the cleared carrier identity. Exact signed
evaluations, evaluations with the actual fixed coefficients, and univariate
polynomial degree replays supplement those exact structural identities.
Thirty deliberate circuit or binding corruptions are rejected. The complete
homogeneous-part argument in Section~\ref{sec:positive}, rather than the
degree of one specialization alone, supplies the global degree proof.

\begin{center}
\begin{tabular}{lr}
\toprule
Principal direct finite certificate & Checked count\\
\midrule
Changed or deleted source rows & 5\\
Recognized affine square roots / Horn product links & 17 / 11\\
Integral coordinate restorations & 222\\
Final integer translations & 13\\
Distinct atoms in the common square budget & 41\\
Independent positive seed indices & 16\\
Collected coefficients in the final quartic & 6,545\\
\bottomrule
\end{tabular}
\end{center}

The supporting semantic regression program checks small factorial candidates,
concrete encoding moduli and accumulators, signed products of consecutive
integers, and the strengthened bound $S_i<V$. It also retains a negative
control showing why arbitrary admissible witnesses cannot be restricted to
squares. Exact test counts are recorded in its receipts.
These tests illustrate the formulas and detect transcription errors;
none is used as a proof of an assertion quantified over all integers.

Supported checkers were run both normally and with Python optimization
enabled, so their correctness checks do not rely on removable \file{assert}
statements. The assembled release was checked from its own directory.
Deterministic regeneration of the principal direct artifacts is separately
recorded. The residual appendix is generated from the complete principal
equation and checked byte-for-byte by
\file{report/generate_residuals.py --check}.

For integration, copy this packet as a new research contribution. Preserve the
supplied source ledgers and the distinction between the new lemmas and the
inherited calibration. The main article's editable source is
\file{report/main.tex}; compile it with a standard TeX Live installation.
The release manifest authenticates the supplied files. Rebuilding receipts or
the PDF can change their bytes; verify the original manifest before rebuilding
if an unchanged-archive check is desired.


\section{Further research questions}
\label{sec:future}

\textbf{Which other witnesses can be chosen canonically?}
The useful criterion is constructive completeness with an explicit order of
choices. Cofinality by itself is insufficient in the presence of residue
conditions. A systematic pass over the remaining free bounds could identify
parameters that can be replaced by existing powers or monotone polynomial
families without introducing fresh coordinates.

\textbf{Can the number of independent exponent certificates be reduced?}
The direct construction shares a Pell frame but retains multiple certified
indices. An identity between genuine requested powers can eliminate a
certificate only if its signs, range bounds, and noncircular recovery order
remain justified. The complete literal source is small enough for such
changes to be measured rather than estimated.

\textbf{Can proof-tree shape improve the accumulator bound?}
The proof above uses the coarse bound $t+1$ on the number of pending
formulas. Depth-first expansion bounds it more sharply in terms of proof-tree
depth. Turning that observation into a smaller polynomial requires encoding
or uniformly controlling that depth without spending more unknowns than are
saved. It may be more useful for exponent size than for variable count.

\textbf{Can the compact positive equations lower both degree and circuit cost?}
The positive companions exchange transport and norm-centering costs for
degree. Their constants, source calibration, and domains should remain fixed
while comparing transformations. A low gate count with very large fixed
coefficients should not be interpreted as a comparable bit-complexity result.

\textbf{What is the best formalization target?}
The new cofinal selection and modulus theorems are short and separable from
the inherited compiler. Formalizing these lemmas and the exact finite
coefficient checker would provide a useful next verification layer. A fully
internal ZFC proof object or an end-to-end Lean/Rocq development remains a
separate task from the ordinary proof delivered here.

\textbf{How should simplicity be scored?}
A reproducible comparison should retain the vector
\[
 (\text{domain},\ \text{variables},\ \text{degree},\
   \text{height},\ \text{monomials},\ \text{live operations})
\]
and report which entries improve. The constructions here do not establish a
global lower bound in any of these measures.

\begin{thebibliography}{9}
\bibitem{repo} Vladimir Reshetnikov, \emph{ProveIt: Diophantine consistency
research}, snapshot
\file{45b399ec51dfbecb891dd51d8a9e4b6b7c361ac6}, 24 October 2026.
\url{https://github.com/VladimirReshetnikov/ProveIt/tree/45b399ec51dfbecb891dd51d8a9e4b6b7c361ac6/Computability/DiophantineConsistency}.
\bibitem{round14} \emph{Diophantine consistency equations with powered
saturation and factorial carriers}, Round 14, 24 October 2026. Incoming
archive \file{ConZFC_Powered_Saturation_Factorial_Carriers_Round14_2026-10-24.zip}
at the snapshot in \cite{repo}. The exact ancestor files and hashes are
retained in this packet's provenance records.
\bibitem{saturated} \emph{ConZFC saturated Pell}, 24 October 2026. Incoming
archive \file{ConZFC_Saturated_Pell_2026-10-24.zip} at the same snapshot.
Its complete source coefficient binding and applicable proof documents are
retained with the positive companions.
\bibitem{structural} \emph{Removing Redundancy from Diophantine Consistency
Equations}, 24 October 2026. Incoming archive
\file{con-zfc-structural-compression-2026-10-24.zip}, same snapshot.
\bibitem{bayer} Jonas Bayer, Marco David, Malte Ha\ss ler, Yuri Matiyasevich,
and Dierk Schleicher, \emph{Diophantine equations over $\Z$: universal bounds
and parallel formalization}, arXiv:2506.20909v3.
\url{https://arxiv.org/html/2506.20909v3}. Used for context and the standard
three-square representation, not as a substituted universal-polynomial
parameter in the delivered direct equation.
\end{thebibliography}

\clearpage
\appendix
% Deterministically generated by generate_residuals.py; do not edit.
% Source SHA256: d10a073dbe4e57421696186eb92fbb4775c15b0b5973a8dd324dec1691c3ab4d
% Source: direct/principal/equation.json (packet-relative).
% Variables: 508; residuals: 308; source terms: 1758.
\section{Complete quadratic residuals of the principal equation}
\label{app:principal-residuals}

The unknowns are $x_0,\ldots,x_{507}\in\mathbb Z$, with the
zero-based coordinate order of \file{direct/principal/equation.json}.
The following list contains every term of all $308$ integral
quadratic residuals. The single principal equation is exactly
\[
  P_{508}(x_0,\ldots,x_{507})
  =\sum_{j=0}^{307}Q_j(x_0,\ldots,x_{507})^2=0.
\]
The symbols $Q_j$ are polynomial abbreviations, not additional
unknowns. Every coefficient is printed below; no term is omitted.

\begingroup
\footnotesize
\allowdisplaybreaks[4]
\setlength{\jot}{2pt}
\begin{align*}
 Q_{0}={}& - 225 - 30\,x_{0} - x_{0}^{2} + x_{4} \\[3pt]
 Q_{1}={}& - 15 - x_{0} + x_{0}\,x_{4} + 15\,x_{4} \\*
 &{} - x_{4}\,x_{190} + x_{4}\,x_{193} + x_{177} - x_{180} \\[3pt]
 Q_{2}={}& - x_{0}\,x_{2} - 16\,x_{2} + x_{9} \\[3pt]
 Q_{3}={}& - 1 + x_{1} - x_{2}\,x_{7} - x_{3} \\*
 &{} - x_{7} \\[3pt]
 Q_{4}={}& - 1 + x_{1} - x_{8} - x_{8}\,x_{9} \\[3pt]
 Q_{5}={}& - 15 - x_{0} - x_{3}^{2} + x_{190} \\*
 &{} - x_{193} \\[3pt]
 Q_{6}={}& - 14 - x_{0} + x_{18} - x_{19}\,x_{21} \\*
 &{} + x_{190} - x_{193} \\[3pt]
 Q_{7}={}& - x_{10}^{2} - 2\,x_{10}\,x_{11} - x_{11}^{2} + x_{22} \\[3pt]
 Q_{8}={}& - 71 - 102\,x_{10} - 36\,x_{10}^{2} - 96\,x_{10}\,x_{11} \\*
 &{} - 24\,x_{10}\,x_{22} - 136\,x_{11} - 64\,x_{11}^{2} - 32\,x_{11}\,x_{22} \\*
 &{} - 34\,x_{22} - 4\,x_{22}^{2} + x_{23} \\[3pt]
 Q_{9}={}& - x_{10}^{2} - 4\,x_{10}\,x_{23} - 4\,x_{23}^{2} + x_{24} \\[3pt]
 Q_{10}={}& - 55 - 90\,x_{10} - 36\,x_{10}^{2} - 192\,x_{10}\,x_{23} \\*
 &{} - 24\,x_{10}\,x_{24} - 240\,x_{23} - 256\,x_{23}^{2} - 64\,x_{23}\,x_{24} \\*
 &{} - 30\,x_{24} - 4\,x_{24}^{2} + x_{25} \\[3pt]
 Q_{11}={}& - x_{11}^{2} - 2\,x_{11}\,x_{12} - x_{12}^{2} + x_{26} \\[3pt]
 Q_{12}={}& - 3\,x_{10} - x_{10}^{2} - 4\,x_{10}\,x_{11} - 4\,x_{10}\,x_{26} \\*
 &{} - 8\,x_{11} - 4\,x_{11}^{2} - 8\,x_{11}\,x_{26} - 8\,x_{26} \\*
 &{} - 4\,x_{26}^{2} + x_{27} \\[3pt]
 Q_{13}={}& - x_{10}^{2} - 2\,x_{10}\,x_{12} - x_{12}^{2} + x_{28} \\[3pt]
 Q_{14}={}& - 150 - 294\,x_{10} - 144\,x_{10}^{2} - 192\,x_{10}\,x_{11} \\*
 &{} - 192\,x_{10}\,x_{12} - 48\,x_{10}\,x_{22} - 48\,x_{10}\,x_{28} - 200\,x_{11} \\*
 &{} - 64\,x_{11}^{2} - 128\,x_{11}\,x_{12} - 32\,x_{11}\,x_{22} - 32\,x_{11}\,x_{28} \\*
 &{} - 192\,x_{12} - 64\,x_{12}^{2} - 32\,x_{12}\,x_{22} - 32\,x_{12}\,x_{28} \\*
 &{} - 50\,x_{22} - 4\,x_{22}^{2} - 8\,x_{22}\,x_{28} - 48\,x_{28} \\*
 &{} - 4\,x_{28}^{2} + x_{29} \\[3pt]
 Q_{15}={}& - 36 - 24\,x_{27} - 4\,x_{27}^{2} - 8\,x_{27}\,x_{29} \\*
 &{} - 24\,x_{29} - 4\,x_{29}^{2} + x_{30} \\[3pt]
 Q_{16}={}& - 4 - 16\,x_{12} - 16\,x_{12}^{2} - 16\,x_{12}\,x_{26} \\*
 &{} - 16\,x_{12}\,x_{28} - 8\,x_{26} - 4\,x_{26}^{2} - 8\,x_{26}\,x_{28} \\*
 &{} - 8\,x_{28} - 4\,x_{28}^{2} + x_{31} \\[3pt]
 Q_{17}={}& - 88 - 114\,x_{10} - 36\,x_{10}^{2} - 96\,x_{10}\,x_{11} \\*
 &{} - 48\,x_{10}\,x_{12} - 24\,x_{10}\,x_{22} - 48\,x_{10}\,x_{28} - 24\,x_{10}\,x_{31} \\*
 &{} - 152\,x_{11} - 64\,x_{11}^{2} - 64\,x_{11}\,x_{12} - 32\,x_{11}\,x_{22} \\*
 &{} - 64\,x_{11}\,x_{28} - 32\,x_{11}\,x_{31} - 72\,x_{12} - 16\,x_{12}^{2} \\*
 &{} - 16\,x_{12}\,x_{22} - 32\,x_{12}\,x_{28} - 16\,x_{12}\,x_{31} - 38\,x_{22} \\*
 &{} - 4\,x_{22}^{2} - 16\,x_{22}\,x_{28} - 8\,x_{22}\,x_{31} - 72\,x_{28} \\*
 &{} - 16\,x_{28}^{2} - 16\,x_{28}\,x_{31} - 36\,x_{31} - 4\,x_{31}^{2} \\*
 &{} + x_{32} \\[3pt]
 Q_{18}={}& - 60 - 15\,x_{10} - x_{10}^{2} - 12\,x_{10}\,x_{11} \\*
 &{} - 16\,x_{10}\,x_{12} - 4\,x_{10}\,x_{26} - 96\,x_{11} - 36\,x_{11}^{2} \\*
 &{} - 96\,x_{11}\,x_{12} - 24\,x_{11}\,x_{26} - 128\,x_{12} - 64\,x_{12}^{2} \\*
 &{} - 32\,x_{12}\,x_{26} - 32\,x_{26} - 4\,x_{26}^{2} + x_{33} \\[3pt]
 Q_{19}={}& - 60 - 96\,x_{10} - 36\,x_{10}^{2} - 12\,x_{10}\,x_{11} \\*
 &{} - 96\,x_{10}\,x_{12} - 24\,x_{10}\,x_{28} - 15\,x_{11} - x_{11}^{2} \\*
 &{} - 16\,x_{11}\,x_{12} - 4\,x_{11}\,x_{28} - 128\,x_{12} - 64\,x_{12}^{2} \\*
 &{} - 32\,x_{12}\,x_{28} - 32\,x_{28} - 4\,x_{28}^{2} + x_{34} \\[3pt]
 Q_{20}={}& - 144 - 48\,x_{33} - 4\,x_{33}^{2} - 8\,x_{33}\,x_{34} \\*
 &{} - 48\,x_{34} - 4\,x_{34}^{2} + x_{35} \\[3pt]
 Q_{21}={}& - 4 - 8\,x_{10} - 4\,x_{10}^{2} - 8\,x_{10}\,x_{12} \\*
 &{} - 8\,x_{10}\,x_{22} - 8\,x_{10}\,x_{26} - 8\,x_{12} - 4\,x_{12}^{2} \\*
 &{} - 8\,x_{12}\,x_{22} - 8\,x_{12}\,x_{26} - 8\,x_{22} - 4\,x_{22}^{2} \\*
 &{} - 8\,x_{22}\,x_{26} - 8\,x_{26} - 4\,x_{26}^{2} + x_{36} \\[3pt]
 Q_{22}={}& - 88 - 152\,x_{10} - 64\,x_{10}^{2} - 160\,x_{10}\,x_{12} \\*
 &{} - 64\,x_{10}\,x_{26} - 32\,x_{10}\,x_{28} - 32\,x_{10}\,x_{36} - 186\,x_{12} \\*
 &{} - 100\,x_{12}^{2} - 80\,x_{12}\,x_{26} - 40\,x_{12}\,x_{28} - 40\,x_{12}\,x_{36} \\*
 &{} - 72\,x_{26} - 16\,x_{26}^{2} - 16\,x_{26}\,x_{28} - 16\,x_{26}\,x_{36} \\*
 &{} - 38\,x_{28} - 4\,x_{28}^{2} - 8\,x_{28}\,x_{36} - 36\,x_{36} \\*
 &{} - 4\,x_{36}^{2} + x_{37} \\[3pt]
 Q_{23}={}& - 55 - 90\,x_{10} - 36\,x_{10}^{2} - 96\,x_{10}\,x_{11} \\*
 &{} - 24\,x_{10}\,x_{22} - 120\,x_{11} - 64\,x_{11}^{2} - 32\,x_{11}\,x_{22} \\*
 &{} - 30\,x_{22} - 4\,x_{22}^{2} + x_{38} \\[3pt]
 Q_{24}={}& - x_{10}^{2} - 4\,x_{10}\,x_{38} - 4\,x_{38}^{2} + x_{39} \\[3pt]
 Q_{25}={}& - 55 - 90\,x_{10} - 36\,x_{10}^{2} - 192\,x_{10}\,x_{38} \\*
 &{} - 24\,x_{10}\,x_{39} - 240\,x_{38} - 256\,x_{38}^{2} - 64\,x_{38}\,x_{39} \\*
 &{} - 30\,x_{39} - 4\,x_{39}^{2} + x_{40} \\[3pt]
 Q_{26}={}& - x_{11}^{2} - 4\,x_{11}\,x_{40} - 4\,x_{40}^{2} + x_{41} \\[3pt]
 Q_{27}={}& - 95 - 158\,x_{10} - 64\,x_{10}^{2} - 128\,x_{10}\,x_{11} \\*
 &{} - 64\,x_{10}\,x_{22} - 160\,x_{11} - 64\,x_{11}^{2} - 64\,x_{11}\,x_{22} \\*
 &{} - 78\,x_{22} - 16\,x_{22}^{2} + x_{42} \\[3pt]
 Q_{28}={}& - 176 - 216\,x_{11} - 64\,x_{11}^{2} - 96\,x_{11}\,x_{12} \\*
 &{} - 32\,x_{11}\,x_{26} - 32\,x_{11}\,x_{42} - 162\,x_{12} - 36\,x_{12}^{2} \\*
 &{} - 24\,x_{12}\,x_{26} - 24\,x_{12}\,x_{42} - 54\,x_{26} - 4\,x_{26}^{2} \\*
 &{} - 8\,x_{26}\,x_{42} - 52\,x_{42} - 4\,x_{42}^{2} + x_{43} \\[3pt]
 Q_{29}={}& - 56 - 120\,x_{10} - 64\,x_{10}^{2} - 96\,x_{10}\,x_{12} \\*
 &{} - 32\,x_{10}\,x_{28} - 32\,x_{10}\,x_{43} - 90\,x_{12} - 36\,x_{12}^{2} \\*
 &{} - 24\,x_{12}\,x_{28} - 24\,x_{12}\,x_{43} - 30\,x_{28} - 4\,x_{28}^{2} \\*
 &{} - 8\,x_{28}\,x_{43} - 28\,x_{43} - 4\,x_{43}^{2} + x_{44} \\[3pt]
 Q_{30}={}& - x_{12}^{2} - 2\,x_{12}\,x_{13} - x_{13}^{2} + x_{45} \\[3pt]
 Q_{31}={}& - x_{10}^{2} - 4\,x_{10}\,x_{12} - 4\,x_{10}\,x_{45} - 4\,x_{12}^{2} \\*
 &{} - 8\,x_{12}\,x_{45} - 4\,x_{45}^{2} + x_{46} \\[3pt]
 Q_{32}={}& - 4\,x_{10}^{2} - 8\,x_{10}\,x_{11} - 8\,x_{10}\,x_{26} - 8\,x_{10}\,x_{28} \\*
 &{} - 4\,x_{11}^{2} - 8\,x_{11}\,x_{26} - 8\,x_{11}\,x_{28} - 4\,x_{26}^{2} \\*
 &{} - 8\,x_{26}\,x_{28} - 4\,x_{28}^{2} + x_{47} \\[3pt]
 Q_{33}={}& - 16\,x_{10}^{2} - 32\,x_{10}\,x_{12} - 32\,x_{10}\,x_{28} - 32\,x_{10}\,x_{45} \\*
 &{} - 16\,x_{10}\,x_{46} - 16\,x_{10}\,x_{47} - 4\,x_{12} - 16\,x_{12}^{2} \\*
 &{} - 32\,x_{12}\,x_{28} - 32\,x_{12}\,x_{45} - 16\,x_{12}\,x_{46} - 16\,x_{12}\,x_{47} \\*
 &{} - 16\,x_{28}^{2} - 32\,x_{28}\,x_{45} - 16\,x_{28}\,x_{46} - 16\,x_{28}\,x_{47} \\*
 &{} - 4\,x_{45} - 16\,x_{45}^{2} - 16\,x_{45}\,x_{46} - 16\,x_{45}\,x_{47} \\*
 &{} - 2\,x_{46} - 4\,x_{46}^{2} - 8\,x_{46}\,x_{47} - 4\,x_{47}^{2} \\*
 &{} + x_{48} \\[3pt]
 Q_{34}={}& - 12 - 39\,x_{12} - 25\,x_{12}^{2} - 40\,x_{12}\,x_{28} \\*
 &{} - 20\,x_{12}\,x_{31} - 32\,x_{28} - 16\,x_{28}^{2} - 16\,x_{28}\,x_{31} \\*
 &{} - 16\,x_{31} - 4\,x_{31}^{2} + x_{49} \\[3pt]
 Q_{35}={}& - 4 - 8\,x_{11} - 4\,x_{11}^{2} - 8\,x_{11}\,x_{26} \\*
 &{} - 8\,x_{26} - 4\,x_{26}^{2} + x_{50} \\[3pt]
 Q_{36}={}& - 4 - 16\,x_{11} - 16\,x_{11}^{2} - 8\,x_{11}\,x_{12} \\*
 &{} - 32\,x_{11}\,x_{26} - 16\,x_{11}\,x_{50} - 4\,x_{12} - x_{12}^{2} \\*
 &{} - 8\,x_{12}\,x_{26} - 4\,x_{12}\,x_{50} - 16\,x_{26} - 16\,x_{26}^{2} \\*
 &{} - 16\,x_{26}\,x_{50} - 8\,x_{50} - 4\,x_{50}^{2} + x_{51} \\[3pt]
 Q_{37}={}& - 3\,x_{11} - x_{11}^{2} - 4\,x_{11}\,x_{12} - 4\,x_{11}\,x_{51} \\*
 &{} - 8\,x_{12} - 4\,x_{12}^{2} - 8\,x_{12}\,x_{51} - 8\,x_{51} \\*
 &{} - 4\,x_{51}^{2} + x_{52} \\[3pt]
 Q_{38}={}& - 49 - 28\,x_{52} - 4\,x_{52}^{2} + x_{53} \\[3pt]
 Q_{39}={}& - 144 - 24\,x_{12} - x_{12}^{2} - 8\,x_{12}\,x_{52} \\*
 &{} - 4\,x_{12}\,x_{53} - 96\,x_{52} - 16\,x_{52}^{2} - 16\,x_{52}\,x_{53} \\*
 &{} - 48\,x_{53} - 4\,x_{53}^{2} + x_{54} \\[3pt]
 Q_{40}={}& - 3\,x_{11} - x_{11}^{2} - 4\,x_{11}\,x_{12} - 4\,x_{11}\,x_{54} \\*
 &{} - 8\,x_{12} - 4\,x_{12}^{2} - 8\,x_{12}\,x_{54} - 8\,x_{54} \\*
 &{} - 4\,x_{54}^{2} + x_{55} \\[3pt]
 Q_{41}={}& - 144 - 48\,x_{49} - 4\,x_{49}^{2} - 8\,x_{49}\,x_{55} \\*
 &{} - 48\,x_{55} - 4\,x_{55}^{2} + x_{56} \\[3pt]
 Q_{42}={}& - 192 - 27\,x_{10} - x_{10}^{2} - 8\,x_{10}\,x_{49} \\*
 &{} - 4\,x_{10}\,x_{56} - 112\,x_{49} - 16\,x_{49}^{2} - 16\,x_{49}\,x_{56} \\*
 &{} - 56\,x_{56} - 4\,x_{56}^{2} + x_{57} \\[3pt]
 Q_{43}={}& - 49 - 28\,x_{57} - 4\,x_{57}^{2} + x_{58} \\[3pt]
 Q_{44}={}& - 144 - 24\,x_{12} - x_{12}^{2} - 8\,x_{12}\,x_{57} \\*
 &{} - 4\,x_{12}\,x_{58} - 96\,x_{57} - 16\,x_{57}^{2} - 16\,x_{57}\,x_{58} \\*
 &{} - 48\,x_{58} - 4\,x_{58}^{2} + x_{59} \\[3pt]
 Q_{45}={}& - 3\,x_{10} - x_{10}^{2} - 4\,x_{10}\,x_{12} - 4\,x_{10}\,x_{59} \\*
 &{} - 8\,x_{12} - 4\,x_{12}^{2} - 8\,x_{12}\,x_{59} - 8\,x_{59} \\*
 &{} - 4\,x_{59}^{2} + x_{60} \\[3pt]
 Q_{46}={}& - 49 - 28\,x_{60} - 4\,x_{60}^{2} + x_{61} \\[3pt]
 Q_{47}={}& - 169 - 52\,x_{11} - 4\,x_{11}^{2} - 8\,x_{11}\,x_{26} \\*
 &{} - 16\,x_{11}\,x_{60} - 8\,x_{11}\,x_{61} - 52\,x_{26} - 4\,x_{26}^{2} \\*
 &{} - 16\,x_{26}\,x_{60} - 8\,x_{26}\,x_{61} - 104\,x_{60} - 16\,x_{60}^{2} \\*
 &{} - 16\,x_{60}\,x_{61} - 52\,x_{61} - 4\,x_{61}^{2} + x_{62} \\[3pt]
 Q_{48}={}& - 12 - 39\,x_{11} - 25\,x_{11}^{2} - 40\,x_{11}\,x_{26} \\*
 &{} - 20\,x_{11}\,x_{62} - 32\,x_{26} - 16\,x_{26}^{2} - 16\,x_{26}\,x_{62} \\*
 &{} - 16\,x_{62} - 4\,x_{62}^{2} + x_{63} \\[3pt]
 Q_{49}={}& - 49 - 28\,x_{55} - 4\,x_{55}^{2} + x_{64} \\[3pt]
 Q_{50}={}& - 324 - 144\,x_{55} - 16\,x_{55}^{2} - 16\,x_{55}\,x_{63} \\*
 &{} - 16\,x_{55}\,x_{64} - 72\,x_{63} - 4\,x_{63}^{2} - 8\,x_{63}\,x_{64} \\*
 &{} - 72\,x_{64} - 4\,x_{64}^{2} + x_{65} \\[3pt]
 Q_{51}={}& - 192 - 27\,x_{12} - x_{12}^{2} - 8\,x_{12}\,x_{63} \\*
 &{} - 4\,x_{12}\,x_{65} - 112\,x_{63} - 16\,x_{63}^{2} - 16\,x_{63}\,x_{65} \\*
 &{} - 56\,x_{65} - 4\,x_{65}^{2} + x_{66} \\[3pt]
 Q_{52}={}& - 5 - 12\,x_{10} - 4\,x_{10}^{2} - 4\,x_{10}\,x_{12} \\*
 &{} - 8\,x_{10}\,x_{22} - 5\,x_{12} - x_{12}^{2} - 4\,x_{12}\,x_{22} \\*
 &{} - 12\,x_{22} - 4\,x_{22}^{2} + x_{67} \\[3pt]
 Q_{53}={}& - 49 - 28\,x_{10} - 4\,x_{10}^{2} - 8\,x_{10}\,x_{22} \\*
 &{} - 8\,x_{10}\,x_{67} - 28\,x_{22} - 4\,x_{22}^{2} - 8\,x_{22}\,x_{67} \\*
 &{} - 28\,x_{67} - 4\,x_{67}^{2} + x_{68} \\[3pt]
 Q_{54}={}& - 66 - 68\,x_{10} - 16\,x_{10}^{2} - 32\,x_{10}\,x_{22} \\*
 &{} - 16\,x_{10}\,x_{66} - 16\,x_{10}\,x_{68} - 68\,x_{22} - 16\,x_{22}^{2} \\*
 &{} - 16\,x_{22}\,x_{66} - 16\,x_{22}\,x_{68} - 32\,x_{66} - 4\,x_{66}^{2} \\*
 &{} - 8\,x_{66}\,x_{68} - 34\,x_{68} - 4\,x_{68}^{2} + x_{69} \\[3pt]
 Q_{55}={}& - 1 - 4\,x_{69} - 4\,x_{69}^{2} + x_{70} \\[3pt]
 Q_{56}={}& - x_{3}\,x_{15} + x_{71} \\[3pt]
 Q_{57}={}& x_{14} - x_{15} - 2\,x_{25}\,x_{71} - x_{72} \\[3pt]
 Q_{58}={}& x_{14} - x_{15} - 4\,x_{27}\,x_{71} - 2\,x_{30}\,x_{71} \\*
 &{} - 12\,x_{71} - x_{73} \\[3pt]
 Q_{59}={}& x_{14} - x_{15} - 2\,x_{32}\,x_{71} - x_{74} \\[3pt]
 Q_{60}={}& x_{14} - x_{15} - 4\,x_{33}\,x_{71} - 2\,x_{35}\,x_{71} \\*
 &{} - 12\,x_{71} - x_{75} \\[3pt]
 Q_{61}={}& x_{14} - x_{15} - 2\,x_{37}\,x_{71} - x_{76} \\[3pt]
 Q_{62}={}& x_{14} - x_{15} - 4\,x_{40}\,x_{71} - 2\,x_{41}\,x_{71} \\*
 &{} - x_{77} \\[3pt]
 Q_{63}={}& x_{14} - x_{15} - 2\,x_{44}\,x_{71} - x_{78} \\[3pt]
 Q_{64}={}& x_{14} - x_{15} - 2\,x_{48}\,x_{71} - x_{79} \\[3pt]
 Q_{65}={}& - x_{10}\,x_{11} + x_{10}\,x_{12} + x_{11}\,x_{12} - x_{12}^{2} \\*
 &{} + x_{80} \\[3pt]
 Q_{66}={}& 1 + 2\,x_{13} + x_{13}^{2} - x_{80}^{2} \\*
 &{} + x_{81} \\[3pt]
 Q_{67}={}& - x_{81}^{2} + x_{82} \\[3pt]
 Q_{68}={}& - 4\,x_{69}\,x_{71} - 2\,x_{70}\,x_{71} + x_{83} \\[3pt]
 Q_{69}={}& - x_{14}^{2} + 2\,x_{14}\,x_{15} + 2\,x_{14}\,x_{83} - x_{15}^{2} \\*
 &{} - 2\,x_{15}\,x_{83} - x_{83}^{2} + x_{84} \\[3pt]
 Q_{70}={}& - x_{3}\,x_{11} + x_{85} \\[3pt]
 Q_{71}={}& - x_{15} - x_{15}\,x_{85} + x_{87}\,x_{88} + x_{88} \\*
 &{} + x_{89} \\[3pt]
 Q_{72}={}& - x_{3}\,x_{10} + x_{86} \\[3pt]
 Q_{73}={}& - 2\,x_{3}\,x_{10} - 2\,x_{3}\,x_{22} + x_{87} \\[3pt]
 Q_{74}={}& - x_{14} - x_{14}\,x_{86} + x_{88} \\[3pt]
 Q_{75}={}& - 6\,x_{3} - 6\,x_{3}\,x_{10} - 8\,x_{3}\,x_{11} - 2\,x_{3}\,x_{22} \\*
 &{} + x_{90} \\[3pt]
 Q_{76}={}& x_{14} + x_{14}\,x_{85} - x_{15} - x_{15}\,x_{90} \\*
 &{} + x_{91} \\[3pt]
 Q_{77}={}& - x_{14}\,x_{72} + x_{15}\,x_{72} + x_{92} \\[3pt]
 Q_{78}={}& - x_{73}\,x_{92} + x_{93} \\[3pt]
 Q_{79}={}& - x_{74}\,x_{93} + x_{94} \\[3pt]
 Q_{80}={}& - x_{75}\,x_{94} + x_{95} \\[3pt]
 Q_{81}={}& - x_{76}\,x_{95} + x_{96} \\[3pt]
 Q_{82}={}& - x_{77}\,x_{96} + x_{97} \\[3pt]
 Q_{83}={}& - x_{78}\,x_{97} + x_{98} \\[3pt]
 Q_{84}={}& - x_{79}\,x_{98} + x_{99} \\[3pt]
 Q_{85}={}& - x_{89}\,x_{99} + x_{100} \\[3pt]
 Q_{86}={}& - x_{91}\,x_{100} + x_{101} \\[3pt]
 Q_{87}={}& - x_{101}^{2} + x_{102} \\[3pt]
 Q_{88}={}& - x_{2} - x_{2}\,x_{18} + x_{103} \\[3pt]
 Q_{89}={}& - x_{16} - x_{16}\,x_{103} + x_{104} \\[3pt]
 Q_{90}={}& - 2\,x_{1}^{2} + 2\,x_{1}\,x_{14} + 2\,x_{1}\,x_{15} + 2\,x_{1}\,x_{104} \\*
 &{} + 2\,x_{1}\,x_{105} - x_{14}^{2} - 2\,x_{14}\,x_{104} - x_{15}^{2} \\*
 &{} - 2\,x_{15}\,x_{105} + x_{20}\,x_{106} - x_{82}\,x_{102} - x_{84}\,x_{102} \\*
 &{} - x_{104}^{2} - x_{105}^{2} \\[3pt]
 Q_{91}={}& - x_{2}\,x_{17} - x_{17} - x_{17}\,x_{103} + x_{105} \\[3pt]
 Q_{92}={}& - 31 - 9\,x_{10} - 9\,x_{11} - 4\,x_{22} \\*
 &{} - 2\,x_{25} - 4\,x_{27} - 2\,x_{30} - 2\,x_{32} \\*
 &{} - 4\,x_{33} - 2\,x_{35} - 2\,x_{37} - 4\,x_{40} \\*
 &{} - 2\,x_{41} - 2\,x_{44} - 2\,x_{48} - 4\,x_{69} \\*
 &{} - 2\,x_{70} + x_{301} - x_{304} \\[3pt]
 Q_{93}={}& 1 + x_{5} - x_{20}\,x_{107} + x_{111} \\[3pt]
 Q_{94}={}& - x_{5}\,x_{110} - x_{110} - x_{110}\,x_{111} + x_{112} \\[3pt]
 Q_{95}={}& - x_{5}\,x_{114} - x_{111}\,x_{114} - x_{114} + x_{115} \\[3pt]
 Q_{96}={}& - x_{20}\,x_{117} + x_{119} \\[3pt]
 Q_{97}={}& - x_{20}\,x_{127} + x_{129} \\[3pt]
 Q_{98}={}& - x_{20}\,x_{133} + x_{135} \\[3pt]
 Q_{99}={}& - x_{20}\,x_{109} + x_{140} \\[3pt]
 Q_{100}={}& - x_{143}\,x_{145} + x_{146} \\[3pt]
 Q_{101}={}& - x_{143}\,x_{146} - x_{146} + x_{147} \\[3pt]
 Q_{102}={}& - 1 + x_{144}^{2} + x_{146}^{2} - x_{147}^{2} \\[3pt]
 Q_{103}={}& - x_{144}\,x_{149} + x_{152} \\[3pt]
 Q_{104}={}& - 1 + x_{148}^{2} + x_{149}^{2} - x_{152}^{2} \\[3pt]
 Q_{105}={}& - 4\,x_{149}\,x_{150} + x_{153} \\[3pt]
 Q_{106}={}& 1 - x_{144} - x_{148}\,x_{151} + x_{153} \\[3pt]
 Q_{107}={}& - x_{144}\,x_{155} + x_{158} \\[3pt]
 Q_{108}={}& - 1 + x_{154}^{2} + x_{155}^{2} - x_{158}^{2} \\[3pt]
 Q_{109}={}& - x_{153}\,x_{157} - x_{157} + x_{159} \\[3pt]
 Q_{110}={}& - 1 + x_{156}^{2} + x_{157}^{2} - x_{159}^{2} \\[3pt]
 Q_{111}={}& - x_{148}\,x_{160} - x_{154} + x_{156} \\[3pt]
 Q_{112}={}& - 15 - x_{0} - 4\,x_{155}\,x_{161} + x_{157} \\[3pt]
 Q_{113}={}& - 197 - 28\,x_{0} - x_{0}^{2} - 2\,x_{0}\,x_{144} \\*
 &{} + 2\,x_{0}\,x_{177} - 2\,x_{0}\,x_{180} - 28\,x_{144} + 2\,x_{144}\,x_{177} \\*
 &{} - 2\,x_{144}\,x_{180} - x_{162} + 28\,x_{177} - x_{177}^{2} \\*
 &{} + 2\,x_{177}\,x_{180} - 28\,x_{180} - x_{180}^{2} \\[3pt]
 Q_{114}={}& x_{0}\,x_{155} + x_{6} - x_{154} + 14\,x_{155} \\*
 &{} - x_{155}\,x_{177} + x_{155}\,x_{180} + x_{158} + x_{162}\,x_{163} \\[3pt]
 Q_{115}={}& - x_{144}\,x_{165} + x_{168} \\[3pt]
 Q_{116}={}& - 1 + x_{164}^{2} + x_{165}^{2} - x_{168}^{2} \\[3pt]
 Q_{117}={}& - x_{153}\,x_{167} - x_{167} + x_{169} \\[3pt]
 Q_{118}={}& - 1 + x_{166}^{2} + x_{167}^{2} - x_{169}^{2} \\[3pt]
 Q_{119}={}& - x_{148}\,x_{170} - x_{164} + x_{166} \\[3pt]
 Q_{120}={}& - 4 - 2\,x_{1} - 2\,x_{2} - 4\,x_{165}\,x_{171} \\*
 &{} + x_{167} - 2\,x_{190} + 2\,x_{193} \\[3pt]
 Q_{121}={}& 1 + x_{19}^{2} - 2\,x_{19}\,x_{144} + x_{172} \\[3pt]
 Q_{122}={}& - x_{19}\,x_{165} + x_{20} - x_{164} + x_{168} \\*
 &{} + x_{172}\,x_{173} \\[3pt]
 Q_{123}={}& - 2 - 2\,x_{6} - x_{6}^{2} - 2\,x_{6}\,x_{20} \\*
 &{} + 2\,x_{6}\,x_{144} - 2\,x_{20} - x_{20}^{2} + 2\,x_{20}\,x_{144} \\*
 &{} + 2\,x_{144} - x_{174} \\[3pt]
 Q_{124}={}& - x_{6}\,x_{165} - x_{20}\,x_{165} + x_{107} - x_{164} \\*
 &{} - x_{165} + x_{168} + x_{174}\,x_{175} \\[3pt]
 Q_{125}={}& - x_{144}\,x_{177} + x_{181} \\[3pt]
 Q_{126}={}& - 1 + x_{176}^{2} + x_{177}^{2} - x_{181}^{2} \\[3pt]
 Q_{127}={}& - x_{153}\,x_{179} - x_{179} + x_{182} \\[3pt]
 Q_{128}={}& - 1 + x_{178}^{2} + x_{179}^{2} - x_{182}^{2} \\[3pt]
 Q_{129}={}& - x_{148}\,x_{183} - x_{176} + x_{178} \\[3pt]
 Q_{130}={}& - x_{177} - 4\,x_{177}\,x_{184} + x_{179} + x_{180} \\[3pt]
 Q_{131}={}& - 5 - 4\,x_{5} - x_{5}^{2} - 2\,x_{5}\,x_{111} \\*
 &{} + 2\,x_{5}\,x_{144} - 4\,x_{111} - x_{111}^{2} + 2\,x_{111}\,x_{144} \\*
 &{} + 4\,x_{144} - x_{185} \\[3pt]
 Q_{132}={}& - x_{5}\,x_{177} + x_{108} - x_{111}\,x_{177} - x_{176} \\*
 &{} - 2\,x_{177} + x_{181} + x_{185}\,x_{186} \\[3pt]
 Q_{133}={}& - 2 - 2\,x_{5} - x_{5}^{2} - 2\,x_{5}\,x_{111} \\*
 &{} + 2\,x_{5}\,x_{144} - 2\,x_{111} - x_{111}^{2} + 2\,x_{111}\,x_{144} \\*
 &{} + 2\,x_{144} - x_{187} \\[3pt]
 Q_{134}={}& - x_{5}\,x_{155} + x_{109} - x_{111}\,x_{155} - x_{154} \\*
 &{} - x_{155} + x_{158} + x_{187}\,x_{188} \\[3pt]
 Q_{135}={}& - x_{144}\,x_{190} + x_{194} \\[3pt]
 Q_{136}={}& - 1 + x_{189}^{2} + x_{190}^{2} - x_{194}^{2} \\[3pt]
 Q_{137}={}& - x_{153}\,x_{192} - x_{192} + x_{195} \\[3pt]
 Q_{138}={}& - 1 + x_{191}^{2} + x_{192}^{2} - x_{195}^{2} \\[3pt]
 Q_{139}={}& - x_{148}\,x_{196} - x_{189} + x_{191} \\[3pt]
 Q_{140}={}& - x_{190} - 4\,x_{190}\,x_{197} + x_{192} + x_{193} \\[3pt]
 Q_{141}={}& - x_{5}\,x_{190} - x_{111}\,x_{190} + x_{113} + x_{185}\,x_{198} \\*
 &{} - x_{189} - 2\,x_{190} + x_{194} \\[3pt]
 Q_{142}={}& - x_{144}\,x_{200} + x_{203} \\[3pt]
 Q_{143}={}& - 1 + x_{199}^{2} + x_{200}^{2} - x_{203}^{2} \\[3pt]
 Q_{144}={}& - x_{153}\,x_{202} - x_{202} + x_{204} \\[3pt]
 Q_{145}={}& - 1 + x_{201}^{2} + x_{202}^{2} - x_{204}^{2} \\[3pt]
 Q_{146}={}& - x_{148}\,x_{205} - x_{199} + x_{201} \\[3pt]
 Q_{147}={}& - x_{10} - 4\,x_{200}\,x_{206} + x_{202} \\[3pt]
 Q_{148}={}& - x_{5}\,x_{200} - x_{111}\,x_{200} + x_{116} + x_{185}\,x_{207} \\*
 &{} - x_{199} - 2\,x_{200} + x_{203} \\[3pt]
 Q_{149}={}& 1 - x_{144}^{2} + x_{208} \\[3pt]
 Q_{150}={}& - x_{5}\,x_{154} - x_{5}\,x_{158} - x_{111}\,x_{154} - x_{111}\,x_{158} \\*
 &{} + x_{117} + x_{144}\,x_{158} - x_{154} - x_{155}\,x_{208} \\*
 &{} - x_{158} + x_{187}\,x_{209} \\[3pt]
 Q_{151}={}& - x_{144}\,x_{211} + x_{214} \\[3pt]
 Q_{152}={}& - 1 + x_{210}^{2} + x_{211}^{2} - x_{214}^{2} \\[3pt]
 Q_{153}={}& - x_{153}\,x_{213} - x_{213} + x_{215} \\[3pt]
 Q_{154}={}& - 1 + x_{212}^{2} + x_{213}^{2} - x_{215}^{2} \\[3pt]
 Q_{155}={}& - x_{148}\,x_{216} - x_{210} + x_{212} \\[3pt]
 Q_{156}={}& - x_{11} - 4\,x_{211}\,x_{217} + x_{213} \\[3pt]
 Q_{157}={}& - x_{5}\,x_{211} - x_{111}\,x_{211} + x_{120} + x_{185}\,x_{218} \\*
 &{} - x_{210} - 2\,x_{211} + x_{214} \\[3pt]
 Q_{158}={}& - x_{144}\,x_{220} + x_{223} \\[3pt]
 Q_{159}={}& - 1 + x_{219}^{2} + x_{220}^{2} - x_{223}^{2} \\[3pt]
 Q_{160}={}& - x_{153}\,x_{222} - x_{222} + x_{224} \\[3pt]
 Q_{161}={}& - 1 + x_{221}^{2} + x_{222}^{2} - x_{224}^{2} \\[3pt]
 Q_{162}={}& - x_{148}\,x_{225} - x_{219} + x_{221} \\[3pt]
 Q_{163}={}& - x_{12} - 4\,x_{220}\,x_{226} + x_{222} \\[3pt]
 Q_{164}={}& - x_{5}\,x_{220} - x_{111}\,x_{220} + x_{122} + x_{185}\,x_{227} \\*
 &{} - x_{219} - 2\,x_{220} + x_{223} \\[3pt]
 Q_{165}={}& - x_{144}\,x_{229} + x_{232} \\[3pt]
 Q_{166}={}& - 1 + x_{228}^{2} + x_{229}^{2} - x_{232}^{2} \\[3pt]
 Q_{167}={}& - x_{153}\,x_{231} - x_{231} + x_{233} \\[3pt]
 Q_{168}={}& - 1 + x_{230}^{2} + x_{231}^{2} - x_{233}^{2} \\[3pt]
 Q_{169}={}& - x_{148}\,x_{234} - x_{228} + x_{230} \\[3pt]
 Q_{170}={}& - x_{13} - 4\,x_{229}\,x_{235} + x_{231} \\[3pt]
 Q_{171}={}& - x_{5}\,x_{229} - x_{111}\,x_{229} + x_{124} + x_{185}\,x_{236} \\*
 &{} - x_{228} - 2\,x_{229} + x_{232} \\[3pt]
 Q_{172}={}& - x_{144}\,x_{238} + x_{241} \\[3pt]
 Q_{173}={}& - 1 + x_{237}^{2} + x_{238}^{2} - x_{241}^{2} \\[3pt]
 Q_{174}={}& - x_{153}\,x_{240} - x_{240} + x_{242} \\[3pt]
 Q_{175}={}& - 1 + x_{239}^{2} + x_{240}^{2} - x_{242}^{2} \\[3pt]
 Q_{176}={}& - x_{148}\,x_{243} - x_{237} + x_{239} \\[3pt]
 Q_{177}={}& - x_{14} - 4\,x_{238}\,x_{244} + x_{240} \\[3pt]
 Q_{178}={}& - x_{5}\,x_{238} - x_{111}\,x_{238} + x_{126} + x_{185}\,x_{245} \\*
 &{} - x_{237} - 2\,x_{238} + x_{241} \\[3pt]
 Q_{179}={}& - x_{144}\,x_{247} + x_{250} \\[3pt]
 Q_{180}={}& - 1 + x_{246}^{2} + x_{247}^{2} - x_{250}^{2} \\[3pt]
 Q_{181}={}& - x_{153}\,x_{249} - x_{249} + x_{251} \\[3pt]
 Q_{182}={}& - 1 + x_{248}^{2} + x_{249}^{2} - x_{251}^{2} \\[3pt]
 Q_{183}={}& - x_{148}\,x_{252} - x_{246} + x_{248} \\[3pt]
 Q_{184}={}& - 1 - x_{2} - 4\,x_{247}\,x_{253} + x_{249} \\[3pt]
 Q_{185}={}& - x_{5}\,x_{247} - x_{111}\,x_{247} + x_{127} + x_{187}\,x_{254} \\*
 &{} - x_{246} - x_{247} + x_{250} \\[3pt]
 Q_{186}={}& - x_{144}\,x_{256} + x_{259} \\[3pt]
 Q_{187}={}& - 1 + x_{255}^{2} + x_{256}^{2} - x_{259}^{2} \\[3pt]
 Q_{188}={}& - x_{153}\,x_{258} - x_{258} + x_{260} \\[3pt]
 Q_{189}={}& - 1 + x_{257}^{2} + x_{258}^{2} - x_{260}^{2} \\[3pt]
 Q_{190}={}& - x_{148}\,x_{261} - x_{255} + x_{257} \\[3pt]
 Q_{191}={}& - x_{15} - 4\,x_{256}\,x_{262} + x_{258} \\[3pt]
 Q_{192}={}& - x_{5}\,x_{256} - x_{111}\,x_{256} + x_{130} + x_{185}\,x_{263} \\*
 &{} - x_{255} - 2\,x_{256} + x_{259} \\[3pt]
 Q_{193}={}& - x_{144}\,x_{265} + x_{268} \\[3pt]
 Q_{194}={}& - 1 + x_{264}^{2} + x_{265}^{2} - x_{268}^{2} \\[3pt]
 Q_{195}={}& - x_{153}\,x_{267} - x_{267} + x_{269} \\[3pt]
 Q_{196}={}& - 1 + x_{266}^{2} + x_{267}^{2} - x_{269}^{2} \\[3pt]
 Q_{197}={}& - x_{148}\,x_{270} - x_{264} + x_{266} \\[3pt]
 Q_{198}={}& - x_{16} - 4\,x_{265}\,x_{271} + x_{267} \\[3pt]
 Q_{199}={}& - x_{5}\,x_{265} - x_{111}\,x_{265} + x_{132} + x_{185}\,x_{272} \\*
 &{} - x_{264} - 2\,x_{265} + x_{268} \\[3pt]
 Q_{200}={}& - x_{144}\,x_{274} + x_{277} \\[3pt]
 Q_{201}={}& - 1 + x_{273}^{2} + x_{274}^{2} - x_{277}^{2} \\[3pt]
 Q_{202}={}& - x_{153}\,x_{276} - x_{276} + x_{278} \\[3pt]
 Q_{203}={}& - 1 + x_{275}^{2} + x_{276}^{2} - x_{278}^{2} \\[3pt]
 Q_{204}={}& - x_{148}\,x_{279} - x_{273} + x_{275} \\[3pt]
 Q_{205}={}& - 1 - x_{1} - 4\,x_{274}\,x_{280} + x_{276} \\[3pt]
 Q_{206}={}& - x_{5}\,x_{274} - x_{111}\,x_{274} + x_{133} + x_{187}\,x_{281} \\*
 &{} - x_{273} - x_{274} + x_{277} \\[3pt]
 Q_{207}={}& - x_{144}\,x_{283} + x_{286} \\[3pt]
 Q_{208}={}& - 1 + x_{282}^{2} + x_{283}^{2} - x_{286}^{2} \\[3pt]
 Q_{209}={}& - x_{153}\,x_{285} - x_{285} + x_{287} \\[3pt]
 Q_{210}={}& - 1 + x_{284}^{2} + x_{285}^{2} - x_{287}^{2} \\[3pt]
 Q_{211}={}& - x_{148}\,x_{288} - x_{282} + x_{284} \\[3pt]
 Q_{212}={}& - x_{17} - 4\,x_{283}\,x_{289} + x_{285} \\[3pt]
 Q_{213}={}& - x_{5}\,x_{283} - x_{111}\,x_{283} + x_{136} + x_{185}\,x_{290} \\*
 &{} - x_{282} - 2\,x_{283} + x_{286} \\[3pt]
 Q_{214}={}& - x_{144}\,x_{292} + x_{295} \\[3pt]
 Q_{215}={}& - 1 + x_{291}^{2} + x_{292}^{2} - x_{295}^{2} \\[3pt]
 Q_{216}={}& - x_{153}\,x_{294} - x_{294} + x_{296} \\[3pt]
 Q_{217}={}& - 1 + x_{293}^{2} + x_{294}^{2} - x_{296}^{2} \\[3pt]
 Q_{218}={}& - x_{148}\,x_{297} - x_{291} + x_{293} \\[3pt]
 Q_{219}={}& - x_{18} - 4\,x_{292}\,x_{298} + x_{294} \\[3pt]
 Q_{220}={}& - x_{5}\,x_{292} - x_{111}\,x_{292} + x_{138} + x_{185}\,x_{299} \\*
 &{} - x_{291} - 2\,x_{292} + x_{295} \\[3pt]
 Q_{221}={}& - x_{144}\,x_{301} + x_{305} \\[3pt]
 Q_{222}={}& - 1 + x_{300}^{2} + x_{301}^{2} - x_{305}^{2} \\[3pt]
 Q_{223}={}& - x_{153}\,x_{303} - x_{303} + x_{306} \\[3pt]
 Q_{224}={}& - 1 + x_{302}^{2} + x_{303}^{2} - x_{306}^{2} \\[3pt]
 Q_{225}={}& - x_{148}\,x_{307} - x_{300} + x_{302} \\[3pt]
 Q_{226}={}& - x_{301} - 4\,x_{301}\,x_{308} + x_{303} + x_{304} \\[3pt]
 Q_{227}={}& - x_{5}\,x_{301} - x_{111}\,x_{301} + x_{141} + x_{185}\,x_{309} \\*
 &{} - x_{300} - 2\,x_{301} + x_{305} \\[3pt]
 Q_{228}={}& - x_{155}\,x_{165} + x_{310} \\[3pt]
 Q_{229}={}& - x_{177}\,x_{310} + x_{311} \\[3pt]
 Q_{230}={}& - x_{190}\,x_{311} + x_{312} \\[3pt]
 Q_{231}={}& - x_{200}\,x_{312} + x_{313} \\[3pt]
 Q_{232}={}& - x_{211}\,x_{313} + x_{314} \\[3pt]
 Q_{233}={}& - x_{220}\,x_{314} + x_{315} \\[3pt]
 Q_{234}={}& - x_{229}\,x_{315} + x_{316} \\[3pt]
 Q_{235}={}& - x_{238}\,x_{316} + x_{317} \\[3pt]
 Q_{236}={}& - x_{247}\,x_{317} + x_{318} \\[3pt]
 Q_{237}={}& - x_{256}\,x_{318} + x_{319} \\[3pt]
 Q_{238}={}& - x_{265}\,x_{319} + x_{320} \\[3pt]
 Q_{239}={}& - x_{274}\,x_{320} + x_{321} \\[3pt]
 Q_{240}={}& - x_{283}\,x_{321} + x_{322} \\[3pt]
 Q_{241}={}& - x_{292}\,x_{322} + x_{323} \\[3pt]
 Q_{242}={}& - x_{301}\,x_{323} + x_{324} \\[3pt]
 Q_{243}={}& - x_{324}^{2} + x_{325} \\[3pt]
 Q_{244}={}& x_{149} - x_{325}\,x_{326} \\[3pt]
 Q_{245}={}& - 703 - 90\,x_{0} - 3\,x_{0}^{2} + 2\,x_{0}\,x_{177} \\*
 &{} - 2\,x_{0}\,x_{180} - 18\,x_{1} - 5\,x_{1}^{2} - 8\,x_{1}\,x_{2} \\*
 &{} - 8\,x_{1}\,x_{190} + 8\,x_{1}\,x_{193} - 18\,x_{2} - 5\,x_{2}^{2} \\*
 &{} - 8\,x_{2}\,x_{190} + 8\,x_{2}\,x_{193} - 6\,x_{5} - 2\,x_{5}^{2} \\*
 &{} - 4\,x_{5}\,x_{111} - 2\,x_{6} - 2\,x_{6}^{2} - 2\,x_{6}\,x_{20} \\*
 &{} - x_{10}^{2} - x_{11}^{2} - x_{12}^{2} - x_{13}^{2} \\*
 &{} - x_{14}^{2} - x_{15}^{2} - x_{16}^{2} - x_{17}^{2} \\*
 &{} - x_{18}^{2} - x_{19}^{2} - 2\,x_{20} - 2\,x_{20}^{2} \\*
 &{} - x_{107}^{2} - x_{108}^{2} - x_{109}^{2} - 6\,x_{111} \\*
 &{} - 2\,x_{111}^{2} - x_{113}^{2} - x_{116}^{2} - x_{117}^{2} \\*
 &{} - x_{120}^{2} - x_{122}^{2} - x_{124}^{2} - x_{126}^{2} \\*
 &{} - x_{127}^{2} - x_{130}^{2} - x_{132}^{2} - x_{133}^{2} \\*
 &{} - x_{136}^{2} - x_{138}^{2} - x_{141}^{2} + x_{143} \\*
 &{} + 28\,x_{177} - 2\,x_{177}^{2} + 4\,x_{177}\,x_{180} - 28\,x_{180} \\*
 &{} - 2\,x_{180}^{2} - 16\,x_{190} - 5\,x_{190}^{2} + 10\,x_{190}\,x_{193} \\*
 &{} + 16\,x_{193} - 5\,x_{193}^{2} - x_{301}^{2} + 2\,x_{301}\,x_{304} \\*
 &{} - x_{304}^{2} \\[3pt]
 Q_{246}={}& - 1 + x_{0} - x_{507}^{2} \\[3pt]
 Q_{247}={}& x_{1} - x_{327}^{2} - x_{328}^{2} - x_{329} \\*
 &{} - x_{329}^{2} \\[3pt]
 Q_{248}={}& - 1 - x_{3}\,x_{5} + x_{5} + x_{6} \\*
 &{} - x_{330}^{2} - x_{331}^{2} - x_{332} - x_{332}^{2} \\[3pt]
 Q_{249}={}& x_{3}\,x_{5} - x_{6} - x_{333}^{2} - x_{334}^{2} \\*
 &{} - x_{335} - x_{335}^{2} \\[3pt]
 Q_{250}={}& x_{2} + x_{6} - x_{6}^{2} \\[3pt]
 Q_{251}={}& - 1 + x_{10} - x_{336}^{2} - x_{337}^{2} \\*
 &{} - x_{338} - x_{338}^{2} \\[3pt]
 Q_{252}={}& - 1 + x_{11} - x_{339}^{2} - x_{340}^{2} \\*
 &{} - x_{341} - x_{341}^{2} \\[3pt]
 Q_{253}={}& - 1 + x_{12} - x_{342}^{2} - x_{343}^{2} \\*
 &{} - x_{344} - x_{344}^{2} \\[3pt]
 Q_{254}={}& - 1 + x_{13} - x_{345}^{2} - x_{346}^{2} \\*
 &{} - x_{347} - x_{347}^{2} \\[3pt]
 Q_{255}={}& - 1 + x_{14} - x_{348}^{2} - x_{349}^{2} \\*
 &{} - x_{350} - x_{350}^{2} \\[3pt]
 Q_{256}={}& - 1 + x_{15} - x_{351}^{2} - x_{352}^{2} \\*
 &{} - x_{353} - x_{353}^{2} \\[3pt]
 Q_{257}={}& - 1 + x_{16} - x_{354}^{2} - x_{355}^{2} \\*
 &{} - x_{356} - x_{356}^{2} \\[3pt]
 Q_{258}={}& - 1 + x_{17} - x_{357}^{2} - x_{358}^{2} \\*
 &{} - x_{359} - x_{359}^{2} \\[3pt]
 Q_{259}={}& - 1 + x_{18} - x_{360}^{2} - x_{361}^{2} \\*
 &{} - x_{362} - x_{362}^{2} \\[3pt]
 Q_{260}={}& x_{19} - x_{363}^{2} - x_{364}^{2} - x_{365} \\*
 &{} - x_{365}^{2} \\[3pt]
 Q_{261}={}& - x_{10} - x_{11} - x_{12} - x_{13} \\*
 &{} - x_{14} - x_{15} - x_{16} - x_{17} \\*
 &{} - x_{18} + 9\,x_{20}^{2} - x_{366}^{2} - x_{367}^{2} \\*
 &{} - x_{368} - x_{368}^{2} \\[3pt]
 Q_{262}={}& - x_{5}\,x_{109} + x_{108} - x_{109}\,x_{112} - x_{369}^{2} \\*
 &{} - x_{370}^{2} - x_{371} - x_{371}^{2} \\[3pt]
 Q_{263}={}& - 1 + x_{5}\,x_{109} - x_{108} + x_{109} \\*
 &{} + x_{109}\,x_{112} - x_{372}^{2} - x_{373}^{2} - x_{374} \\*
 &{} - x_{374}^{2} \\[3pt]
 Q_{264}={}& x_{111} - x_{375}^{2} - x_{376}^{2} - x_{377} \\*
 &{} - x_{377}^{2} \\[3pt]
 Q_{265}={}& - x_{19}\,x_{109} - x_{109}\,x_{115} + x_{113} - x_{378}^{2} \\*
 &{} - x_{379}^{2} - x_{380} - x_{380}^{2} \\[3pt]
 Q_{266}={}& - 1 + x_{19}\,x_{109} + x_{109} + x_{109}\,x_{115} \\*
 &{} - x_{113} - x_{381}^{2} - x_{382}^{2} - x_{383} \\*
 &{} - x_{383}^{2} \\[3pt]
 Q_{267}={}& x_{5} - x_{19} + x_{111} - x_{384}^{2} \\*
 &{} - x_{385}^{2} - x_{386} - x_{386}^{2} \\[3pt]
 Q_{268}={}& x_{116} - x_{118}\,x_{119} - x_{387}^{2} - x_{388}^{2} \\*
 &{} - x_{389} - x_{389}^{2} \\[3pt]
 Q_{269}={}& - 1 - x_{116} + x_{117} + x_{118}\,x_{119} \\*
 &{} - x_{390}^{2} - x_{391}^{2} - x_{392} - x_{392}^{2} \\[3pt]
 Q_{270}={}& - x_{119}\,x_{121} + x_{120} - x_{393}^{2} - x_{394}^{2} \\*
 &{} - x_{395} - x_{395}^{2} \\[3pt]
 Q_{271}={}& - 1 + x_{117} + x_{119}\,x_{121} - x_{120} \\*
 &{} - x_{396}^{2} - x_{397}^{2} - x_{398} - x_{398}^{2} \\[3pt]
 Q_{272}={}& - x_{119}\,x_{123} + x_{122} - x_{399}^{2} - x_{400}^{2} \\*
 &{} - x_{401} - x_{401}^{2} \\[3pt]
 Q_{273}={}& - 1 + x_{117} + x_{119}\,x_{123} - x_{122} \\*
 &{} - x_{402}^{2} - x_{403}^{2} - x_{404} - x_{404}^{2} \\[3pt]
 Q_{274}={}& - x_{119}\,x_{125} + x_{124} - x_{405}^{2} - x_{406}^{2} \\*
 &{} - x_{407} - x_{407}^{2} \\[3pt]
 Q_{275}={}& - 1 + x_{117} + x_{119}\,x_{125} - x_{124} \\*
 &{} - x_{408}^{2} - x_{409}^{2} - x_{410} - x_{410}^{2} \\[3pt]
 Q_{276}={}& x_{126} - x_{128}\,x_{129} - x_{411}^{2} - x_{412}^{2} \\*
 &{} - x_{413} - x_{413}^{2} \\[3pt]
 Q_{277}={}& - 1 - x_{126} + x_{127} + x_{128}\,x_{129} \\*
 &{} - x_{414}^{2} - x_{415}^{2} - x_{416} - x_{416}^{2} \\[3pt]
 Q_{278}={}& - x_{129}\,x_{131} + x_{130} - x_{417}^{2} - x_{418}^{2} \\*
 &{} - x_{419} - x_{419}^{2} \\[3pt]
 Q_{279}={}& - 1 + x_{127} + x_{129}\,x_{131} - x_{130} \\*
 &{} - x_{420}^{2} - x_{421}^{2} - x_{422} - x_{422}^{2} \\[3pt]
 Q_{280}={}& x_{132} - x_{134}\,x_{135} - x_{423}^{2} - x_{424}^{2} \\*
 &{} - x_{425} - x_{425}^{2} \\[3pt]
 Q_{281}={}& - 1 - x_{132} + x_{133} + x_{134}\,x_{135} \\*
 &{} - x_{426}^{2} - x_{427}^{2} - x_{428} - x_{428}^{2} \\[3pt]
 Q_{282}={}& - x_{135}\,x_{137} + x_{136} - x_{429}^{2} - x_{430}^{2} \\*
 &{} - x_{431} - x_{431}^{2} \\[3pt]
 Q_{283}={}& - 1 + x_{133} + x_{135}\,x_{137} - x_{136} \\*
 &{} - x_{432}^{2} - x_{433}^{2} - x_{434} - x_{434}^{2} \\[3pt]
 Q_{284}={}& x_{138} - x_{139}\,x_{140} - x_{435}^{2} - x_{436}^{2} \\*
 &{} - x_{437} - x_{437}^{2} \\[3pt]
 Q_{285}={}& - 1 + x_{109} - x_{138} + x_{139}\,x_{140} \\*
 &{} - x_{438}^{2} - x_{439}^{2} - x_{440} - x_{440}^{2} \\[3pt]
 Q_{286}={}& - x_{119}\,x_{142} + x_{141} - x_{441}^{2} - x_{442}^{2} \\*
 &{} - x_{443} - x_{443}^{2} \\[3pt]
 Q_{287}={}& - 1 + x_{117} + x_{119}\,x_{142} - x_{141} \\*
 &{} - x_{444}^{2} - x_{445}^{2} - x_{446} - x_{446}^{2} \\[3pt]
 Q_{288}={}& x_{144} - x_{447}^{2} - x_{448}^{2} - x_{449} \\*
 &{} - x_{449}^{2} \\[3pt]
 Q_{289}={}& - 1 + x_{145} - x_{450}^{2} - x_{451}^{2} \\*
 &{} - x_{452} - x_{452}^{2} \\[3pt]
 Q_{290}={}& - 1 + x_{150} - x_{453}^{2} - x_{454}^{2} \\*
 &{} - x_{455} - x_{455}^{2} \\[3pt]
 Q_{291}={}& - 15 - x_{0} + x_{155} - x_{456}^{2} \\*
 &{} - x_{457}^{2} - x_{458} - x_{458}^{2} \\[3pt]
 Q_{292}={}& - 4 - 2\,x_{1} - 2\,x_{2} + x_{165} \\*
 &{} - 2\,x_{190} + 2\,x_{193} - x_{459}^{2} - x_{460}^{2} \\*
 &{} - x_{461} - x_{461}^{2} \\[3pt]
 Q_{293}={}& x_{180} - x_{462}^{2} - x_{463}^{2} - x_{464} \\*
 &{} - x_{464}^{2} \\[3pt]
 Q_{294}={}& x_{193} - x_{465}^{2} - x_{466}^{2} - x_{467} \\*
 &{} - x_{467}^{2} \\[3pt]
 Q_{295}={}& - x_{10} + x_{200} - x_{468}^{2} - x_{469}^{2} \\*
 &{} - x_{470} - x_{470}^{2} \\[3pt]
 Q_{296}={}& - x_{11} + x_{211} - x_{471}^{2} - x_{472}^{2} \\*
 &{} - x_{473} - x_{473}^{2} \\[3pt]
 Q_{297}={}& - x_{12} + x_{220} - x_{474}^{2} - x_{475}^{2} \\*
 &{} - x_{476} - x_{476}^{2} \\[3pt]
 Q_{298}={}& - x_{13} + x_{229} - x_{477}^{2} - x_{478}^{2} \\*
 &{} - x_{479} - x_{479}^{2} \\[3pt]
 Q_{299}={}& - x_{14} + x_{238} - x_{480}^{2} - x_{481}^{2} \\*
 &{} - x_{482} - x_{482}^{2} \\[3pt]
 Q_{300}={}& - 1 - x_{2} + x_{247} - x_{483}^{2} \\*
 &{} - x_{484}^{2} - x_{485} - x_{485}^{2} \\[3pt]
 Q_{301}={}& - x_{15} + x_{256} - x_{486}^{2} - x_{487}^{2} \\*
 &{} - x_{488} - x_{488}^{2} \\[3pt]
 Q_{302}={}& - x_{16} + x_{265} - x_{489}^{2} - x_{490}^{2} \\*
 &{} - x_{491} - x_{491}^{2} \\[3pt]
 Q_{303}={}& - 1 - x_{1} + x_{274} - x_{492}^{2} \\*
 &{} - x_{493}^{2} - x_{494} - x_{494}^{2} \\[3pt]
 Q_{304}={}& - x_{17} + x_{283} - x_{495}^{2} - x_{496}^{2} \\*
 &{} - x_{497} - x_{497}^{2} \\[3pt]
 Q_{305}={}& - x_{18} + x_{292} - x_{498}^{2} - x_{499}^{2} \\*
 &{} - x_{500} - x_{500}^{2} \\[3pt]
 Q_{306}={}& x_{304} - x_{501}^{2} - x_{502}^{2} - x_{503} \\*
 &{} - x_{503}^{2} \\[3pt]
 Q_{307}={}& - 1 + x_{326} - x_{504}^{2} - x_{505}^{2} \\*
 &{} - x_{506} - x_{506}^{2}
\end{align*}
\endgroup


\end{document}
